\documentclass[a4paper,11pt]{article}
\pdfoutput=1
\usepackage{jheppub}   % JHEP class file (SISSA Medialab, v1.1227), copied into this directory
\usepackage[T1]{fontenc}
\usepackage{amsmath,amssymb,amsfonts,mathtools}
\usepackage{braket}
\usepackage{booktabs}
\usepackage{xcolor}
\usepackage{tikz}
\usetikzlibrary{calc}
\usetikzlibrary{decorations.pathmorphing}
\usetikzlibrary{decorations.markings}
\usetikzlibrary{decorations.text}
\usetikzlibrary{decorations.pathreplacing}
\usepackage{tikzit}
\input{stiles.tikzstyles}

% ---------------------------------------------------------------------
% Notation (identical to arXiv:2607.18373 and to the calculation files)
% ---------------------------------------------------------------------
\def\mf#1{\mathbb{#1}}          % coefficients of Omega: \mf{N},\mf{A},\mf{B},\mf{C}
\def\v#1{\mathbf{#1}}
\newcommand{\Kv}{\mathcal{K}}    % Weizsaecker-Williams field  K^i(x)= x^i/x^2
\newcommand{\MS}{\overline{\rm MS}}
\newcommand{\Tr}{{\rm Tr}}
\newcommand{\NLOfirst}{\mathrm{NLO}}
\newcommand{\Lam}{\Lambda}
% upper cut-off of the soft modes is denoted \vee, as in 2607.18373 and 1610.03453 (\vee = \Lambda e^{\delta Y})

\title{Single inclusive gluon production at mid rapidity at next-to-leading order in the Color Glass Condensate}

\author{AUTHORS}
\affiliation{AFFILIATIONS}
% \emailAdd{}

\abstract{We compute the single inclusive gluon spectrum in the collision of a dilute projectile with a dense target at next-to-leading order in the Color Glass Condensate framework, for a gluon produced at mid rapidity, in pure Yang--Mills theory. The calculation is performed in light-cone perturbation theory: the projectile is described by its soft gluon light-cone wave function, generated from the valence colour charge density $\rho$ by a unitary dressing operator $\Omega$ known to the required order, and the target by eikonal adjoint Wilson lines. The valence charges are treated as non-commuting operators throughout. We organise the $O(g^4)$ cross section into interference terms with one gluon in the final state and into contributions with two gluons in the final state, one of which is integrated over. All terms enhanced by the logarithm of the longitudinal cut-off ratio exponentiate as expected: the logarithm splits at the longitudinal momentum $k^+$ of the produced gluon, and the coefficient of $\ln(\vee/k^+)$ is exactly the leading-order BFKL evolution of the projectile colour charge correlator, while the coefficient of $\ln(k^+/\Lambda)$ is exactly the JIMWLK evolution of the target Wilson-line operator of the leading-order cross section. The terms quartic and (in the symmetrised sense) cubic in $\rho$ that accompany these logarithms cancel. The remaining contributions consist of a collinear (DGLAP) term associated with the fragmentation of the produced gluon, a running-coupling term, and finite terms, which we give in colour-singlet form.
% NOTE (draft): the last sentence must be revised once sections \ref{sec:dglap_rc} and \ref{sec:finite} are filled in.
}

\begin{document}
\maketitle
\flushbottom

% ======================================================================
\section{Introduction}\label{sec:intro}
% ======================================================================

At high energy the wave function of a hadron contains a large number of soft gluons, and in a sufficiently dense system their density saturates~\cite{Gribov:1983ivg,Mueller:1985wy}. The Color Glass Condensate (CGC) is the effective theory of this regime~\cite{McLerran:1993ni,McLerran:1993ka,Iancu:2003xm,Weigert:2005us,Gelis:2010nm,Kovchegov:2012mbw}: the fast (valence) degrees of freedom of a hadron act as static colour sources, characterised by a colour charge density $\rho^a(\v x)$, for the soft gluon field, and the dependence of observables on the rapidity interval over which the soft modes are integrated out is governed by the BFKL equation~\cite{Kuraev:1977fs,Balitsky:1978ic} in the dilute regime and by the non-linear Balitsky--Kovchegov (BK) and JIMWLK equations~\cite{Balitsky:1995ub,Kovchegov:1999yj,Kovchegov:1999ua,Jalilian-Marian:1997qno,Jalilian-Marian:1997jx,Jalilian-Marian:1997ubg,Kovner:2000pt,Weigert:2000gi,Iancu:2000hn,Iancu:2001ad,Ferreiro:2001qy} in the dense regime.

The simplest observable that is sensitive both to the projectile and to the target wave functions is the single inclusive gluon spectrum in the collision of a dilute projectile with a dense target. At leading order (LO), for a gluon produced at mid rapidity, i.e.\ far in rapidity from both the projectile and the target, the spectrum is given by the formula of Kovchegov and Tuchin~\cite{Kovchegov:2001sc,Kovchegov:1998bi}: the gluon is emitted from the valence charges of the projectile in the Weizs\"acker--Williams field, and the emission amplitude, together with the colour rotation of the emitting charge, scatters eikonally on the target. The cross section is a product of the projectile colour charge correlator and of a target operator built from two adjoint Wilson lines; in the limit of a dilute target it reduces to the $k_T$-factorised form with the Lipatov vertex~\cite{Gribov:1983ivg,Kharzeev:2003wz,Blaizot:2004wu}. The leading logarithmic (LL) rapidity evolution of this result was obtained by Kovner and Lublinsky~\cite{Kovner:2006wr} (see also~\cite{Kovner:2006ge}) and, within the Schwinger--Keldysh formulation, by Gelis, Lappi and Venugopalan~\cite{Gelis:2008rw}: the projectile is to be evolved from its own rapidity down to the rapidity of the produced gluon, and the target from the rapidity of the produced gluon to its own rapidity. For a dilute projectile the first evolution is linear (BFKL) and acts on the projectile charge correlator, while the second is the JIMWLK evolution of the target Wilson lines.

To make quantitative use of this result, the next-to-leading order (NLO) corrections, i.e.\ the terms of relative order $\alpha_s$ that are not enhanced by a rapidity logarithm, are required. The ingredients of such a calculation have become available over the last decade. The evolution equations themselves are known at NLO: the NLO BK equation~\cite{Balitsky:2007feb} and the NLO JIMWLK Hamiltonian~\cite{Balitsky:2013fea,Kovner:2013ona,Kovner:2014lca,Grabovsky:2013mba}, together with the running-coupling corrections to the LO kernels~\cite{Balitsky:2006wa,Kovchegov:2006vj,Kovchegov:2006wf}. In particular, Lublinsky and Mulian~\cite{Lublinsky:2016meo} computed the soft gluon light-cone wave function of a dilute projectile to order $g^3$ and derived the NLO JIMWLK Hamiltonian from it. For particle production, the NLO corrections to single inclusive hadron production at forward rapidity were computed in the hybrid formalism~\cite{Dumitru:2005gt,Altinoluk:2011qy} by Chirilli, Xiao and Yuan~\cite{Chirilli:2011km,Chirilli:2012jd}, and the structure of the rapidity subtraction and of the kinematic constraints in this framework was analysed in~\cite{Stasto:2013cha,Altinoluk:2014eka,Ducloue:2016shw}. Running-coupling corrections to inclusive gluon production at mid rapidity were obtained by Horowitz and Kovchegov~\cite{Horowitz:2010yg}, and the relation between $\beta_0$-dependent terms of NLO calculations and the running of the coupling was clarified in~\cite{Kovner:2023vsy}. The general structure of the NLO correction to inclusive spectra in the CGC was discussed by Gelis and Laidet~\cite{Gelis:2012ri}. Finally, corrections of higher order in the projectile colour charge (``first saturation corrections'' in the projectile) were computed in the classical approximation in~\cite{Chirilli:2015tea,Li:2021zmf,Li:2021ntm}, and their contribution to single inclusive production, odd in $\v k$, was analysed in~\cite{McLerran:2016snu,Kovchegov:2018jun}. A complete NLO calculation of the single inclusive gluon spectrum at mid rapidity has, to our knowledge, not been presented.

In a previous paper~\cite{Radhakrishnan:2026soft} the soft gluon wave function of a dilute projectile and the associated unitary evolution (dressing) operator $\Omega$ were constructed up to order $g^2$, in pure Yang--Mills theory, within light-cone perturbation theory (LCPT)~\cite{Lepage:1980fj,Brodsky:1997de}. The NLO cross section requires in addition the coefficient of $\Omega$ that creates one gluon at order $g^3$; its ingredients, the $O(g^3)$ one-gluon components of the light-cone wave function, were computed by Lublinsky and Mulian~\cite{Lublinsky:2016meo}, and we recompute them here with the longitudinal cut-offs and the operator ordering required by the cross section. The operator $\Omega$ was written as a normal-ordered series in soft gluon creation and annihilation operators with coefficients that are functionals of the non-commuting valence charge density $\rho$, and the coefficients were fixed by matching to LCPT and by unitarity. In the present paper we use $\Omega$ to compute the single inclusive gluon spectrum at mid rapidity at NLO. The observable is the expectation value of the number operator of dressed gluons in the scattered, dressed projectile state~\cite{Kovner:2006wr}, so that the full calculation reduces to an algebraic evaluation of products of the coefficients of $\Omega$ and of Wilson lines, followed by longitudinal and transverse integrations.

Our main results are the following.
\begin{itemize}
\item The $O(g^4)$ cross section naturally decomposes into three groups of terms: interference of the LO amplitude with the $O(g^3)$ one-gluon amplitude (Group~I); the squares of the $O(g^2)$ two-gluon amplitudes, with one of the gluons integrated over (Group~II); and the interference terms between different two-gluon amplitudes (Group~III). We list all products of coefficients that contribute, with their longitudinal integration ranges and their order in $\rho$ (Section~\ref{sec:observable}).
\item All terms enhanced by $\log(\vee/\Lambda)$, where $\Lambda<p^+<\vee$ is the longitudinal momentum range of the soft modes, are the LL evolution of the LO result. The logarithm splits at the longitudinal momentum $k^+$ of the produced gluon into $\log(\vee/k^+)+\log(k^+/\Lambda)$; the coefficient of $\log(\vee/k^+)$ is exactly the BFKL evolution of the projectile charge correlator inserted into the LO formula, and the coefficient of $\log(k^+/\Lambda)$ is exactly the JIMWLK Hamiltonian acting on the target operator of the LO formula. We give both terms for an arbitrary (ordered) projectile correlator and in colour-singlet form. The contributions quartic in $\rho$, all of which carry $\log(\vee/\Lambda)$, cancel identically for commuting charges, and the totally symmetrised cubic terms cancel as well; with non-commuting charges, their reordering produces the two-$\rho$ terms that complete the evolution (Section~\ref{sec:rapidity}).
\item The remaining terms are a collinear term, associated with the fragmentation of the produced gluon and governed by the gluon--gluon splitting function, a running-coupling term, and finite terms. We give an account of all poles $1/\epsilon$: the ultraviolet poles of the virtual corrections to the produced gluon cancel among themselves, the coupling is renormalised by the normalisation of the dressed produced gluon, and the collinear poles combine into the full $\MS$ splitting function $\mathcal P_{gg}$ and are absorbed into the fragmentation function. The running-coupling term is the $\frac{\beta_0}{2}\delta(1-\xi)$ part of the collinear term; it sets the scale of the LO coupling (Section~\ref{sec:dglap_rc}).
\item We give the finite terms for an arbitrary ordered charge correlator and, averaged over the projectile, in colour-singlet form. The double logarithms $\log^2(k^+/\Lambda)$ of individual rows cancel in the sum. The soft remainders of the plus prescriptions cancel in pairs of a virtual and a real term with the same colour structure and kernel, without any condition on the target at large distance; in particular we show which terms cancel the non-integrable logarithm of one of the virtual rows. After the average over the projectile the terms cubic in $\rho$ do not vanish: their symmetrised part, odd in $\v k$, is proportional to the $d^{abc}$ (odderon) three-point function of the projectile, kept general, and the commutators of the ordered charges leave contact terms, even in $\v k$, which are present for any projectile; both are given row by row (Section~\ref{sec:finite}).
\end{itemize}
The result can be written schematically as
\begin{align}\label{eq:schematic}
\frac{d^3N}{d^3k} = \frac{d^3N_{\rm LO}}{d^3k}
+ \log\frac{\vee}{k^+}\,\frac{d^3N_{\rm BFKL}}{d^3k}
+ \log\frac{k^+}{\Lambda}\,\frac{d^3N_{\rm JIMWLK}}{d^3k}
+ \frac{d^3N_{\rm DGLAP}}{d^3k}
+ \frac{d^3N_{\rm RC}}{d^3k}
+ \frac{d^3N_{\rm finite}}{d^3k}\,,
\end{align}
where every term on the right-hand side except the first is of relative order $\alpha_s$.

The paper is organised as follows. In Section~\ref{sec:setup} we summarise the setup: light-cone conventions, the Hamiltonian of the soft modes in the eikonal approximation, LCPT with operator-valued matrix elements, and the expansion of $\Omega$ with its coefficients in coordinate space. In Section~\ref{sec:observable} we define the observable, rederive the LO cross section and organise the NLO contributions. Section~\ref{sec:rapidity} is devoted to the rapidity logarithms and their identification with the BFKL and JIMWLK evolution. The collinear and running-coupling terms are discussed in Section~\ref{sec:dglap_rc}, the finite terms in Section~\ref{sec:finite}, and limits and comparisons with the literature in Section~\ref{sec:checks}. We conclude in Section~\ref{sec:conclusions}. The appendices collect the Hamiltonian and its matrix elements (Appendix~\ref{app:hamiltonian}), the light-cone wave functions from which the coefficients of $\Omega$ are obtained (Appendix~\ref{app:wavefunctions}), the unitarity relations among the coefficients (Appendix~\ref{app:unitarity}), the Fourier transforms (Appendix~\ref{app:fourier}), the prescriptions and distributions used in the longitudinal and transverse integrals (Appendix~\ref{app:distributions}) and the complete list of the finite terms for an arbitrary ordered charge correlator (Appendix~\ref{app:finite}).

% ======================================================================
\section{Setup: soft gluon wave function and the evolution operator}\label{sec:setup}
% ======================================================================

In this section we summarise the framework of~\cite{Radhakrishnan:2026soft}, on which the present calculation is built, and collect the coefficients of the evolution operator $\Omega$ that enter the single inclusive cross section. The coefficients of order $g$ and $g^2$ were derived in~\cite{Radhakrishnan:2026soft}; the coefficient of order $g^3$ that creates one gluon, $\mf A^{(3)}_2$, is obtained here. The derivations of the light-cone wave functions from which the coefficients are obtained are summarised in Appendix~\ref{app:wavefunctions}.

\subsection{Light-cone conventions}\label{sec:conventions}

We use light-cone coordinates $x^\pm=(x^0\pm x^3)/\sqrt2$ and denote transverse vectors by $\v x=(x^1,x^2)$; where no confusion is possible the transverse vectors are written without boldface, and $\int_{x}\equiv\int d^2x$. The projectile moves in the $+$ direction, $x^+$ plays the role of time, and $p^+>0$ is the longitudinal momentum of a physical particle. The light-cone energy of an on-shell gluon is $p^-=\v p^2/2p^+$, and we write
\begin{align}
&E_g(p)=\frac{\v p^2}{2p^+}\,,\qquad E_{gg}(p,q)=\frac{\v p^2}{2p^+}+\frac{\v q^2}{2q^+}\,,\notag\\ &d^3p\equiv dp^+d^2p\,,\qquad \delta^{(3)}(p-q)\equiv\delta(p^+-q^+)\,\delta^{(2)}(\v p-\v q)\,.
\end{align}
We work in the light-cone gauge $A^+=0$, in which the transverse components $A_i^a$, $i=1,2$, are the dynamical fields. The gluon creation and annihilation operators satisfy
\begin{align}\label{eq:aadagger}
\big[a^a_i(k),a^{\dagger b}_j(p)\big]=(2\pi)^3\,\delta^{ab}\,\delta_{ij}\,\delta^{(3)}(k-p)\,.
\end{align}
For the Fourier transforms in the transverse plane we use the conventions
\begin{align}\label{eq:FTconventions}
&a(k^+,\v k)=\int d^2z\,e^{-i\v k\cdot\v z}\,a(k^+,\v z)\,,\qquad
a^\dagger(k^+,\v k)=\int d^2z\,e^{i\v k\cdot\v z}\,a^\dagger(k^+,\v z)\,,\notag\\
&f(\v k)=\int d^2x\,e^{i\v k\cdot\v x}\,f(\v x)\,,
\end{align}
so that $\big[a^a_i(k^+,\v z),a^{\dagger b}_j(p^+,\v z')\big]=2\pi\,\delta^{ab}\delta_{ij}\,\delta(k^+-p^+)\,\delta^{(2)}(\v z-\v z')$. One- and two-gluon states are normalised as
\begin{align}\label{eq:states}
\ket{g^a_i(k^+,\v z)}=\frac{a^{\dagger a}_i(k^+,\v z)}{(2\pi)^{1/2}}\ket{0}\,,\qquad
\ket{g^a_i(k^+,\v z)\,g^b_j(p^+,\v z')}=\frac{a^{\dagger a}_i(k^+,\v z)\,a^{\dagger b}_j(p^+,\v z')}{2\pi}\ket{0}\,,
\end{align}
with $\braket{g^b_j(p^+,\v z')|g^a_i(k^+,\v z)}=\delta^{ab}\delta_{ij}\,\delta(k^+-p^+)\,\delta^{(2)}(\v z-\v z')$; in momentum space $\ket{g^a_i(k)}=(2\pi)^{-3/2}a^{\dagger a}_i(k)\ket0$ (Appendix~\ref{app:states}).

Throughout the paper we use the Weizs\"acker--Williams field of a point charge and the transverse kernels built from it,
\begin{align}\label{eq:kernels}
\Kv^i(\v x)\equiv\frac{x^i}{\v x^2}\,,\qquad
K(\v x,\v y,\v z)\equiv\Kv(\v x-\v z)\cdot\Kv(\v y-\v z)=\frac{(\v x-\v z)\cdot(\v y-\v z)}{(\v x-\v z)^2(\v y-\v z)^2}\,,
\end{align}
so that $K(\v x,\v x,\v z)=1/(\v x-\v z)^2$ and $K(\v x,\v x,\v z)-2K(\v x,\v y,\v z)+K(\v y,\v y,\v z)=(\v x-\v y)^2/[(\v x-\v z)^2(\v y-\v z)^2]$ is the dipole kernel.

\subsection{Soft and valence modes; the eikonal Hamiltonian}\label{sec:hamiltonian}

Following the CGC logic~\cite{McLerran:1993ni,McLerran:1993ka,Jalilian-Marian:1997jx,Kovner:2005en}, the gluon modes of the projectile are divided according to their longitudinal momentum. The valence modes have $k^+>\vee$, the soft modes have
\begin{align}
\Lambda<k^+<\vee\,,\qquad \vee\equiv\Lambda\,e^{\delta Y}\,,
\end{align}
and the modes with $k^+<\Lambda$ do not belong to the projectile wave function; their effect is part of the target. The rapidity interval $\delta Y=\ln(\vee/\Lambda)$ is the interval over which the soft modes are integrated; the notation is that of~\cite{Lublinsky:2016meo,Radhakrishnan:2026soft}. Since the valence modes are much faster than the soft ones, they act on the soft modes only through their colour charge density
\begin{align}\label{eq:rho}
\rho^a(\v x)\,,\qquad \big[\rho^a(\v x),\rho^b(\v y)\big]=i f^{abc}\,\rho^c(\v x)\,\delta^{(2)}(\v x-\v y)\,,\qquad
\big[\rho^a(\v k),\rho^b(\v p)\big]=i f^{abc}\,\rho^c(\v k+\v p)\,,
\end{align}
which is an operator on the valence Hilbert space; its microscopic definition in terms of the valence gluon operators is given in~\cite{Radhakrishnan:2026soft}. We keep the charges non-commuting throughout. This is necessary: products of $\rho$'s appear in the cross section in a definite order, determined by the time ordering of the light-cone perturbation theory, and the commutators~\eqref{eq:rho} reduce the number of $\rho$'s by one. These reductions produce some of the logarithmic terms of Section~\ref{sec:rapidity}.

After the decomposition of the gauge field into valence and soft components, all terms suppressed by powers of $k^+_{\rm soft}/k^+_{\rm valence}$ are dropped (eikonal approximation). The light-cone Yang--Mills Hamiltonian ($N_f=0$) then takes the form~\cite{Lublinsky:2016meo,Radhakrishnan:2026soft}
\begin{align}\label{eq:Hdecomp}
&H=H_0+H_g+H_{ggg}+H_{gg\text{-inst}}+H_{gggg}+\dots,\notag\\
&H_0=\int_\Lambda^\vee\frac{dk^+}{2\pi}\int\frac{d^2k}{(2\pi)^2}\,\frac{\v k^2}{2k^+}\,a^{\dagger a}_i(k^+,\v k)\,a^a_i(k^+,\v k)\,,
\end{align}
where the dots stand for terms which do not contribute at the order considered. The coupling of the soft gluons to the valence charges is
\begin{align}\label{eq:Hg}
H_g=\int_\Lambda^\vee\frac{dk^+}{2\pi}\int\frac{d^2k}{(2\pi)^2}\,\frac{g\,k^i}{\sqrt2\,|k^+|^{3/2}}\Big[a^{\dagger a}_i(k^+,\v k)\,\rho^a(-\v k)+a^a_i(k^+,\v k)\,\rho^a(\v k)\Big]\,.
\end{align}
It describes the emission and absorption of a soft gluon by the valence charges. $H_{ggg}$ is the triple-gluon vertex among the soft modes, $H_{gg\text{-inst}}$ the instantaneous (Coulomb-like) interaction between two soft gluons and the valence charge density, and $H_{gggg}$ the four-gluon interaction of the soft modes. Their explicit forms and the matrix elements used in this paper are listed in Appendix~\ref{app:hamiltonian}.

\subsection{Light-cone perturbation theory with operator-valued matrix elements}\label{sec:lcpt}

The soft gluon wave function of the projectile is obtained in time-independent (old-fashioned) perturbation theory in $H-H_0$ on the light cone~\cite{Lepage:1980fj,Brodsky:1997de}. Since the matrix elements of $H_g$ and $H_{gg\text{-inst}}$ are operators on the valence space, the perturbative series is written with the matrix elements in the order of the light-cone time ordering: the vertex acting later stands to the left. With $E_0=0$ for the soft vacuum, the first terms of the expansion of the dressed state in terms of free Fock states are
\begin{align}\label{eq:lcpt}
\ket{\Psi}=\ &\ket0-\sum_n\ket n\frac{V_{n0}}{E_n}+\sum_{n,m}\ket n\frac{V_{nm}V_{m0}}{E_nE_m}
-\sum_{n,m,l}\ket n\frac{V_{nm}V_{ml}V_{l0}}{E_nE_mE_l}
+\sum_{n,m}\ket n\frac{V_{n0}\,V_{0m}V_{m0}}{E_n^2E_m}
\notag\\
&-\frac12\sum_m\frac{V_{0m}V_{m0}}{E_m^2}\ket0
+\frac12\sum_{n,m}\ket n\frac{V_{n0}\,V_{0m}V_{m0}}{E_nE_m^2}+\dots,
\end{align}
where $V_{nm}=\braket{n|H-H_0|m}$ and the sums run over Fock states of the soft sector. The first line is the Rayleigh--Schr\"odinger series with intermediate normalisation; the second line implements the canonical normalisation of the state and is needed at orders $g^2$ (vacuum term) and $g^3$ (one-gluon term). In the last term of each line the single-emission matrix element $V_{n0}$ stands to the left of the norm factor $V_{0m}V_{m0}$; this ordering follows from the perturbative construction of the unitary operator that diagonalises $H$ with operator-valued couplings, and we have verified it against an explicit construction of $\Omega$ with finite-dimensional matrices representing the charges. Energy denominators that vanish inside an integration region are understood as principal values (Appendix~\ref{app:distributions}).

\subsection{The evolution operator}\label{sec:omega}

The dressed (physical) states of the projectile are obtained from the free Fock states by a unitary operator $\Omega$, $\Omega^\dagger\Omega=1$. Since the soft modes are treated in the eikonal approximation, $\Omega$ can be written as a normal-ordered series in the soft gluon operators, with coefficients that are functionals of $\rho$~\cite{Radhakrishnan:2026soft}:
\begin{align}\label{eq:OmegaExp}
\Omega = \ &\mf{N}
+\int\frac{d^3p}{(2\pi)^3}\Big[\mf{A}^b_{1j}(p)\,a^b_j(p)+\mf{A}^b_{2j}(p)\,a^{\dagger b}_j(p)\Big]
\notag\\&
+\int\frac{d^3p}{(2\pi)^3}\frac{d^3q}{(2\pi)^3}\Big[\mf{B}^{cb}_{1kj}(p,q)\,a^b_j(p)a^c_k(q)+\mf{B}^{cb}_{2kj}(p,q)\,a^{\dagger b}_j(p)a^{\dagger c}_k(q)
\notag\\&\hspace{3.4cm}+\mf{B}^{cb}_{3kj}(p,q)\,a^{\dagger c}_k(q)a^b_j(p)\Big]
\notag\\&
+\int\frac{d^3p}{(2\pi)^3}\frac{d^3q}{(2\pi)^3}\Big[\mf{C}^{dcb}_{1lkj}(p,q)\,a^{\dagger d}_l(q)a^{\dagger c}_k(p-q)a^b_j(p)
\notag\\&\hspace{3.4cm}+\mf{C}^{dcb}_{2lkj}(p,q)\,a^{\dagger b}_j(p)a^c_k(p-q)a^d_l(q)\Big]+\dots
\end{align}
The dots denote terms with more gluon operators, which do not contribute to the single inclusive cross section at $O(g^4)$. The coefficients have definite orders in $g$: $\mf N=1+\mf N_2$ with $\mf N_2=O(g^2)$; $\mf A_{1,2}=\mf A^{(1)}_{1,2}+\mf A^{(3)}_{1,2}$ with $\mf A^{(n)}=O(g^n)$; $\mf B_{1,2,3}=O(g^2)$; $\mf C_{1,2}=O(g)$. The coefficients that multiply only creation operators, $\mf A_2$ and $\mf B_2$, together with $\mf N$, are fixed by the action of $\Omega$ on the soft vacuum, i.e.\ by the light-cone wave function of the projectile; $\mf B_3$ and $\mf C_1$ are fixed by the action of $\Omega$ on a one-gluon state. The remaining coefficients follow from unitarity, imposed order by order in $g$ (Appendix~\ref{app:unitarity}):
\begin{align}\label{eq:unitarity_short}
&\mf A^{(1)\dagger b}_{1j}(p)=-\mf A^{(1)b}_{2j}(p)\,,\qquad
\mf C^{\dagger dcb}_{2lkj}(p,q)=-\mf C^{dcb}_{1lkj}(p,q)\,,\notag\\
&\mf N=1-\frac12\int\frac{d^3p}{(2\pi)^3}\,\mf A^{(1)b}_{1j}(p)\,\mf A^{(1)\dagger b}_{1j}(p)\,,
\end{align}
together with the relations for $\mf A^{(3)}_1$ and $\mf B_1$ given in Eqs.~\eqref{eq:A31unit} and~\eqref{eq:B1unit}. In the relation for $\mf A^{(3)}_1$ the operator character of the coefficients matters: $\mf N_2$ does not commute with $\mf A^{(1)}$, since $[\rho^a(\v x)\rho^a(\v y),\rho^b(\v w)]\neq0$ for $\v x\ne\v y$.

In transverse coordinate space, after the unitarity relations~\eqref{eq:unitarity_short} are used, the part of $\Omega$ that contributes is
\begin{align}\label{eq:OmegaX}
\Omega = \ &\mf N
+\int\frac{dp^+}{2\pi}\int_x\Big[\mf A^{(1)b}_{1j}(p^+,\v x)\,a^b_j(p^+,\v x)-\mf A^{\dagger(1)b}_{1j}(p^+,\v x)\,a^{\dagger b}_j(p^+,\v x)
\notag\\&\hspace{2.6cm}
+\mf A^{(3)b}_{1j}(p^+,\v x)\,a^b_j(p^+,\v x)+\mf A^{(3)b}_{2j}(p^+,-\v x)\,a^{\dagger b}_j(p^+,\v x)\Big]
\notag\\&
+\int\frac{dp^+}{2\pi}\frac{dq^+}{2\pi}\int_{x,y}\Big[\mf B^{cb}_{1kj}(p^+,\v x;q^+,\v y)\,a^b_j(p^+,\v x)a^c_k(q^+,\v y)
\notag\\&\hspace{3.2cm}
+\mf B^{cb}_{2kj}(p^+,-\v x;q^+,-\v y)\,a^{\dagger b}_j(p^+,\v x)a^{\dagger c}_k(q^+,\v y)
\notag\\&\hspace{3.2cm}
+\mf B^{cb}_{3kj}(p^+,\v x;q^+,-\v y)\,a^{\dagger c}_k(q^+,\v y)a^b_j(p^+,\v x)\Big]
\notag\\&
+\int\frac{dp^+}{2\pi}\frac{dq^+}{2\pi}\int_{x,y,z}\Big[\mf C^{dcb}_{1lkj}(p^+,\v z-\v y;q^+,\v y-\v x)
\notag\\&\hspace{3.2cm}\times a^{\dagger d}_l(q^+,\v x)a^{\dagger c}_k(p^+-q^+,\v y)a^b_j(p^+,\v z)
\notag\\&\hspace{3.2cm}
-\mf C^{\dagger dcb}_{1lkj}(p^+,\v z-\v y;q^+,\v y-\v x)
\notag\\&\hspace{3.2cm}\times a^{\dagger b}_j(p^+,\v z)a^c_k(p^+-q^+,\v y)a^d_l(q^+,\v x)\Big]+\dots
\end{align}
The transverse arguments of the coefficients are fixed by the Fourier conventions~\eqref{eq:FTconventions}.

\subsection{The coefficients}\label{sec:coefficients}

We now list the coefficients in transverse coordinate space. They are obtained by comparing $\Omega\ket0$ and $\Omega\ket{g^a_i(k^+,\v z)}$ with the LCPT wave functions of Appendix~\ref{app:wavefunctions}, sector by sector in Fock space; the factors of $2\pi$ come from the normalisation~\eqref{eq:states}. The momentum-space expressions and the intermediate steps are given in Appendix~\ref{app:wavefunctions} (for the coefficients of order $g$ and $g^2$ also in~\cite{Radhakrishnan:2026soft}).

\paragraph{Normalisation.} The coefficient of the identity is fixed by the norm of the LO wave function,
\begin{align}\label{eq:N2}
\mf N=1+\mf N_2\,,\qquad
\mf N_2&=-\frac{g^2}{8\pi^3}\,\ln\frac{\vee}{\Lambda}\int_{x,y,z}K(\v x,\v y,\v z)\,\rho^a(\v x)\rho^a(\v y)
\notag\\&=-\frac{g^2}{8\pi^3}\,\ln\frac{\vee}{\Lambda}\int d^2p\,\frac{\rho^a(-\v p)\rho^a(\v p)}{\v p^2}\,.
\end{align}
$\mf N_2$ is the (virtual) probability of not emitting a soft gluon anywhere in the interval $\Lambda<p^+<\vee$, and it carries the full logarithm $\ln(\vee/\Lambda)$.

\paragraph{Single-gluon emission at $O(g)$.} The Weizs\"acker--Williams amplitude for the emission of a soft gluon at $\v z$ is
\begin{align}\label{eq:A1}
\mf A^{(1)a}_{2i}(k^+,-\v z)=-\mf A^{(1)\dagger a}_{1i}(k^+,\v z)=\frac{ig}{\sqrt2\,\pi\sqrt{k^+}}\int_x\Kv^i(\v x-\v z)\,\rho^a(\v x)\,.
\end{align}
It corresponds to the LO wave function $\ket{\Psi^{\rm LO}}=\frac{ig}{2\pi^{3/2}}\int_\Lambda^\vee\frac{dk^+}{\sqrt{k^+}}\int_{x,z}\Kv^i(\v x-\v z)\,\rho^a(\v x)\ket{g^a_i(k^+,\v z)}$, Eq.~\eqref{eq:PsiLO}. In what follows we use the shorthand $\mf A\equiv\mf A^{(1)}_2$ when no confusion can arise, and a bar denotes the rotation of all charges by the target, $\bar{\mf A}^a(\v x)=\mf A^a(\v x)\big|_{\rho(\v u)\to U(\v u)\rho(\v u)}$, see Section~\ref{sec:observable}.

\paragraph{Single-gluon emission at $O(g^3)$.} The coefficient $\mf A^{(3)}_2$ contains terms with one, two and three charges. In momentum space the one-$\rho$ part reads, for $\Lambda\to0$ wherever this is possible,
\begin{align}\label{eq:A3onerho}
\Big[\mf A^{(3)a}_{2i}(k^+,\v k)\Big]_{\rho}=\ &\frac{g^3N_c\,\rho^a(-\v k)\,k^i}{8\sqrt{2k^+}\,\pi^2\,\v k^2}\Bigg\{\Big[-\frac2\epsilon+\ln\frac{\v k^2}{\mu^2_{\MS}}\Big]
\notag\\&\hspace{1.5cm}\times\Big[\frac{11}{3}-2\ln\frac{k^+}{\Lambda}+\ln\Big(1-\frac{k^+}{\vee}\Big)\Big]
\notag\\&\hspace{1.5cm}
+\frac{\pi^2}{3}-\frac{67}{9}+\ln^2\Big(1-\frac{k^+}{\vee}\Big)+\ln\Big(1-\frac{k^+}{\vee}\Big)\ln\frac{\vee}{k^+}\Bigg\},
\end{align}
where the transverse integrals are performed in $d=2-\epsilon$ dimensions in the $\MS$ scheme. It collects the one-loop self-energy of the emitted gluon, the vertex corrections in which the emitted gluon splits and recombines, and the one-$\rho$ (commutator) parts of the diagrams with two valence vertices (Figures~\ref{fig:vac_g3_onerho} and~\ref{fig:vac_g3_tworho}). The pole multiplies the combination $\frac{11}{3}-2\ln(k^+/\Lambda)$, whose first term is $\beta_0/N_c$ for pure glue, $\beta_0=\frac{11}{3}N_c$. In coordinate space,
\begin{align}\label{eq:A3onerhoX}
\Big[\mf A^{(3)a}_{2i}(k^+,-\v z)\Big]_{\rho}=\ &-\int_x\frac{ig^3N_c\,\rho^a(\v x)\,\Kv^i(\v x-\v z)}{16\sqrt2\,\pi^3\sqrt{k^+}}\Bigg\{\Big[-\frac2\epsilon-2\gamma_E-\ln\frac{(\v x-\v z)^2\mu^2_{\MS}}{4}\Big]
\notag\\&\hspace{1.5cm}\times\Big[\frac{11}{3}-2\ln\frac{k^+}{\Lambda}+\ln\Big(1-\frac{k^+}{\vee}\Big)\Big]
\notag\\&\hspace{1.5cm}
-\frac{67}{9}+\frac{\pi^2}{3}+\ln^2\Big(1-\frac{k^+}{\vee}\Big)+\ln\Big(1-\frac{k^+}{\vee}\Big)\ln\frac{\vee}{k^+}\Bigg\}.
\end{align}
The two-$\rho$ part is antisymmetric in colour, $\propto f^{abc}\{\rho^b,\rho^c\}$, and is conveniently written in a mixed representation,
\begin{align}\label{eq:A3tworho}
\Big[\mf A^{(3)a}_{2i}(k^+,-\v z)\Big]_{\rho\rho}=\ &\int d^2p\,d^2k\int_{x,y}e^{i\v k\cdot(\v z-\v y)}e^{i\v p\cdot(\v y-\v x)}\,\frac{ig^3f^{abc}\{\rho^b(\v x),\rho^c(\v y)\}}{32\sqrt{2k^+}\,\pi^5\,\v p^2\,\v k^2}
\notag\\&\times\Bigg[\ln\frac{\Lambda}{k^+}\,\frac{\v k^2p^i-(\v k\cdot\v p)\,k^i}{(\v k-\v p)^2}
-2k^i\ln\frac{\vee}{\Lambda}
-2k^i\ln\frac{k^+\v p^2}{\vee\,\v k^2+k^+\v p^2}
\notag\\&\quad
+\frac{(2\v p^2-2\v k\cdot\v p+\v k^2)\,k^i+\v k^2p^i}{(\v k-\v p)^2}\ln\Big[1+\Big(\frac{\vee}{k^+}-1\Big)\frac{\v k^2}{(\v k-\v p)^2}\Big]
\notag\\&
-\frac{\v p^2k^i+\v k^2p^i}{(\v k-\v p)^2}\ln\frac{\vee-k^+}{\Lambda}-k^i\ln\frac{\vee}{k^+}\Bigg],
\end{align}
where, inside the integral, $\v k$ and $\v p$ are integration variables. The three-$\rho$ part originates from the third-order terms of Eq.~\eqref{eq:lcpt} in which a gluon is emitted and reabsorbed by the valence charges before or after the emission of the observed gluon (Figure~\ref{fig:vac_g3_threerho}). The terms proportional to $1/\Lambda$ and $1/\vee$ cancel exactly, and the remainder is
\begin{align}\label{eq:A3threerho}
\Big[\mf A^{(3)a}_{2i}(k^+,-\v z)\Big]_{\rho\rho\rho}&=-\ln\frac{\vee}{\Lambda}\int_{w,x,y,u}\frac{ig^3\,\rho^a(\v w)\,\rho^b(\v x)\,\rho^b(\v y)\,K(\v x,\v y,\v u)\,\Kv^i(\v w-\v z)}{8\sqrt2\,\pi^4\sqrt{k^+}}
\notag\\&=\mf A^{(1)a}_{2i}(k^+,-\v z)\,\mf N_2\,.
\end{align}
The operator ordering in Eq.~\eqref{eq:A3threerho}, with $\rho^a(\v w)$ to the left of the norm factor, is dictated by Eq.~\eqref{eq:lcpt}. It is important in what follows: since $\mf N_2\mf A^{(1)}\neq\mf A^{(1)}\mf N_2$, the opposite ordering would differ by $-[\mf N_2,\mf A^{(1)}]$, which is a two-$\rho$ term equal and opposite to the term $-2k^i\ln(\vee/\Lambda)$ of Eq.~\eqref{eq:A3tworho}.

The one-$\rho$, two-$\rho$ and three-$\rho$ parts of the wave function from which $\mf A^{(3)}_2$ is obtained agree with the NLO light-cone wave function of Lublinsky and Mulian~\cite{Lublinsky:2016meo} (eqs.~3.97--3.102 there), including the coefficient $11/3$, the constant $67/9-\pi^2/3$, the dependence on the longitudinal cut-offs and the three-$\rho$ term; the two-$\rho$ parts agree after the antisymmetrisation enforced by $f^{abc}\{\rho^b,\rho^c\}$. The coordinate-space forms~\eqref{eq:A3onerhoX}--\eqref{eq:A3threerho} go beyond the published material.

\paragraph{Two-gluon emission.} The coefficient $\mf B_2$ gives the amplitude for two soft gluons, with longitudinal momenta $k^+$ and $p^+$, at transverse positions $\v w$ and $\v z$:
\begin{align}\label{eq:B2}
\mf B^{bc}_{2ji}(k^+,-\v w;p^+,-\v z)=\ &\int_x\frac{ig^2f^{cba}\,\rho^a(\v x)\,\sqrt{p^+k^+}}{4\pi^2\big[p^+(\v x-\v z)^2+k^+(\v x-\v w)^2\big](p^++k^+)}
\notag\\&\times\Bigg[\frac{\delta^{ij}\big[(\v x-\v z)^2-(\v x-\v w)^2\big]}{2(\v z-\v w)^2}
\notag\\&\quad
+\frac{p^++k^+}{k^+}\Big(\frac{(x-z)^j(z-w)^i}{(\v z-\v w)^2}-\frac{(x-z)^j(x-w)^i}{2(\v x-\v w)^2}\Big)
\notag\\&\quad
+\frac{p^++k^+}{p^+}\Big(\frac{(x-w)^i(z-w)^j}{(\v z-\v w)^2}+\frac{(x-z)^j(x-w)^i}{2(\v x-\v z)^2}\Big)\Bigg]
\notag\\&\quad\times\Theta(\vee-k^+-p^+)
\notag\\&
-\int_{x,y}\frac{g^2\,\{\rho^c(\v y),\rho^b(\v x)\}\,\Kv^i(\v y-\v w)\,\Kv^j(\v x-\v z)}{8\pi^2\sqrt{p^+k^+}}
\notag\\&\quad\times\Theta(\vee-k^+)\,\Theta(\vee-p^+)\,,
\end{align}
with $k^+,p^+>\Lambda$. The first term, linear in $\rho$, arises from the emission of a gluon of momentum $k^++p^+$ which splits into two (triple-gluon vertex), from the instantaneous emission of the pair, and from the commutator part of the emission of two gluons by two charges (Figure~\ref{fig:vac_gg}); the parent gluon is a soft mode, hence $k^++p^+<\vee$. The second term is the independent emission of the two gluons by two valence charges; each gluon is separately a soft mode, and there is no condition on $k^++p^+$.
% CHOICE FOR THE AUTHORS: the support of the commutator part of the two-emission diagram (contained in the first term) is taken as k^+ + p^+ < vee, as for the splitting diagrams; by the kinematics of that diagram it could also be the full square Lambda<k^+,p^+<vee.

\paragraph{Scattering of a soft gluon on the valence charges.} The coefficient $\mf B_3$ describes the transition of a soft gluon of momentum $k^+$ at $\v z$ into a soft gluon of momentum $p^+$ at $\v u$ through its interaction with the valence charges: instantaneous interaction, emission of a second gluon followed by absorption of the first, and absorption of the gluon followed by re-emission (Figure~\ref{fig:g_g}). For $p^+>k^+$ the valence charges transfer longitudinal momentum to the gluon, for $p^+<k^+$ the gluon transfers momentum to the valence charges. With $\Theta_>\equiv\Theta(p^+-k^+-\Lambda)$ and $\Theta_<\equiv\Theta(k^+-\Lambda-p^+)$,
\begin{align}\label{eq:B3}
\mf B^{ba}_{3ji}(k^+,\v z;p^+,-\v u)=\ &\frac{ig^2f^{abc}}{(2\pi)^2}\int_x\frac{\sqrt{p^+k^+}\,(p^++k^+)}{(k^+-p^+)^2}\,\frac{\rho^c(\v x)\,\delta_{ij}}{p^+(\v x-\v u)^2-k^+(\v x-\v z)^2}
\notag\\&\qquad\times\big[\Theta_>+\Theta_<\big]
\notag\\&
+\frac{ig^2f^{abc}}{2\pi^2k^+}\int_x\frac{\sqrt{k^+p^+}}{k^+-p^+}
\notag\\&\qquad\times\frac{\big(k^+(x-z)^m-p^+(x-u)^m\big)\big((x-z)^l-(x-u)^l\big)}{k^+(\v x-\v z)^2-p^+(\v x-\v u)^2}\,\rho^c(\v x)
\notag\\&\qquad\times\Big(\delta_{il}\delta_{jm}-\frac{k^+}{p^+}\delta_{jl}\delta_{im}-\frac{k^+}{k^+-p^+}\delta_{ij}\delta_{lm}\Big)
\notag\\&\qquad\times\Bigg[-\frac{\Theta_>}{(\v u-\v z)^2}-\frac{p^+k^+\,\Theta_<}{\big(k^+(\v x-\v z)-p^+(\v x-\v u)\big)^2}\Bigg]
\notag\\&
-\frac{g^2}{4\pi^2\sqrt{p^+k^+}}\int_{x,y}\Kv^i(\v x-\v z)\,\Kv^j(\v y-\v u)\,\{\rho^a(\v x),\rho^b(\v y)\}
\notag\\&
-\frac{ig^2f^{abc}}{4\pi^2\sqrt{p^+k^+}}\int_x\Kv^i(\v x-\v z)\,\Kv^j(\v x-\v u)
\notag\\&\qquad\times\frac{k^+(\v x-\v z)^2+p^+(\v x-\v u)^2}{p^+(\v x-\v u)^2-k^+(\v x-\v z)^2}\,\rho^c(\v x)\,.
\end{align}
The double pole $1/(k^+-p^+)^2$ at the boundary $|p^+-k^+|=\Lambda$ cancels between the instantaneous and the non-instantaneous terms of the first two lines, so that no term proportional to $k^+/\Lambda$ arises. The vanishing denominators $p^+(\v x-\v u)^2=k^+(\v x-\v z)^2$ are principal values.

\paragraph{Gluon splitting.} The coefficient $\mf C_1$ is independent of $\rho$; it is the light-cone amplitude for the splitting of a soft gluon $(k^+,\v z)$ into two soft gluons of momenta $p^+$ and $k^+-p^+$ (Figure~\ref{fig:g_gg}):
\begin{align}\label{eq:C1}
\mf C^{bca}_{1jki}(k^+,\v z;p^+,\v x)=\ &-\frac{gf^{abc}}{2\pi\sqrt{2k^+}}\,\frac{\sqrt{p^+(k^+-p^+)}}{k^+}\,\delta^{(2)}\Big(\v z+\frac{p^+}{k^+}\v x\Big)\,\Kv^m(\v x)
\notag\\&\times\Big[\delta_{im}\delta_{jk}-\frac{k^+}{p^+}\delta_{jm}\delta_{ik}-\frac{k^+}{k^+-p^+}\delta_{km}\delta_{ij}\Big]\Theta(k^+-p^+-\Lambda)\,.
\end{align}
The $\delta$-function expresses the conservation of the transverse centre of mass $p^+\v x_1+(k^+-p^+)\v x_2=k^+\v x_0$ in the splitting.

\paragraph{Coefficients fixed by unitarity.} $\mf A^{(3)}_1$ and $\mf B_1$, which multiply annihilation operators, are not independent; they are given in terms of the coefficients above by Eqs.~\eqref{eq:A31unit} and~\eqref{eq:B1unit}. In the cross section they appear only through these relations.

% ======================================================================
\section{The observable and its expansion}\label{sec:observable}
% ======================================================================

\subsection{Single inclusive gluon spectrum}\label{sec:observable_def}

The projectile, before the collision, is described by its valence state $\ket{\Psi_{\rm v}}$ times the dressing of the soft modes, $\Omega\ket{0}$, where $\ket0$ is the vacuum of the soft modes. The target is described by the eikonal scattering matrix $S$. At high energy the interaction of a soft gluon and of the valence charges with the target is a colour rotation by the adjoint Wilson line $U(\v x)$ at fixed transverse position:
\begin{align}\label{eq:Srotation}
S^\dagger\rho^a(\v x)S&=U^{ab}(\v x)\,\rho^b(\v x)\,,\notag\\
S^\dagger a^a_i(k^+,\v x)S&=U^{ab}(\v x)\,a^b_i(k^+,\v x)\,,\qquad
S^\dagger a^{\dagger a}_i(k^+,\v x)S=U^{ab}(\v x)\,a^{\dagger b}_i(k^+,\v x)\,.
\end{align}
After the collision the state is $S\,\Omega\ket0$. A gluon observed in the final state is a dressed (physical) gluon; the number of dressed gluons is counted by the bare number operator in the basis rotated back by $\Omega^\dagger$~\cite{Kovner:2006wr}. The single inclusive spectrum is therefore
\begin{align}\label{eq:observable}
\frac{d^3N}{d^3k}=\frac{1}{(2\pi)^3}\int_{w,w'}e^{-i\v k\cdot(\v w-\v w')}\,
\Big\langle\bra0\,\Omega^\dagger S^\dagger\Omega\;a^{\dagger a}_i(k^+,\v w')\,a^a_i(k^+,\v w)\;\Omega^\dagger S\,\Omega\ket0\Big\rangle\,,
\end{align}
where $\langle\dots\rangle$ denotes the average over the valence state, i.e.\ over the projectile charge configurations, taken at the end. The factor $(2\pi)^{-3}$ and the phase follow from $a^{\dagger a}_i(k)a^a_i(k)=\int_{w,w'}e^{-i\v k\cdot(\v w-\v w')}a^{\dagger a}_i(k^+,\v w')a^a_i(k^+,\v w)$ and the normalisation~\eqref{eq:aadagger}. The amplitude stands at $\v w$ and the conjugate amplitude at $\v w'$; we use this convention, $e^{-i\v k\cdot(\v w-\v w')}$ with the conjugated factor at the primed points, throughout the paper.

Inserting $SS^\dagger=1$ and defining $\bar\Omega\equiv S^\dagger\Omega S$, which is $\Omega$ with $\rho\to U\rho$ and $a\to Ua$ according to Eq.~\eqref{eq:Srotation}, we obtain
\begin{align}
\frac{d^3N}{d^3k}=\frac{1}{(2\pi)^3}\int_{w,w'}e^{-i\v k\cdot(\v w-\v w')}\,\Big\langle\bra0\,\Omega^\dagger\bar\Omega\;U^{ab'}(\v w')a^{\dagger b'}_i(k^+,\v w')\,U^{ab}(\v w)a^b_i(k^+,\v w)\;\bar\Omega^\dagger\Omega\ket0\Big\rangle\,.
\end{align}
A bar on any coefficient of $\Omega$ denotes, here and below, the replacement $\rho^a(\v u)\to U^{ab}(\v u)\rho^b(\v u)$ of every charge in it. Between $a(k^+,\v w)$ and $\bar\Omega^\dagger\Omega\ket0$ we insert a complete set of soft states. The vacuum does not contribute, and states with three or more gluons contribute at order $g^6$. Hence
\begin{align}\label{eq:N_amplitudes}
\frac{d^3N}{d^3k}=\frac{1}{(2\pi)^3}\Big\langle\mathcal A^{a\dagger}_i(k)\,\mathcal A^a_i(k)+\int dq^+\!\int_{v}\mathcal A^{ab\dagger}_{ij}(k;q^+,\v v)\,\mathcal A^{ab}_{ij}(k;q^+,\v v)\Big\rangle\,,
\end{align}
with the one-gluon and two-gluon amplitudes
\begin{align}\label{eq:amplitudes_def}
\mathcal A^a_i(k)&=\sqrt{2\pi}\int_{w}e^{-i\v k\cdot\v w}\,U^{ab}(\v w)\,\bra{g^b_i(k^+,\v w)}\bar\Omega^\dagger\Omega\ket0\,,\notag\\
\mathcal A^{ab}_{ij}(k;q^+,\v v)&=\sqrt{2\pi}\int_{w}e^{-i\v k\cdot\v w}\,U^{ac}(\v w)\,\bra{g^b_j(q^+,\v v)\,g^c_i(k^+,\v w)}\bar\Omega^\dagger\Omega\ket0\,.
\end{align}
The colour index $b$ of the second gluon in $\mathcal A^{ab}_{ij}$ is that of the bare state; since the second gluon is integrated over, its rotation by $U(\v v)$ drops out of the sum. The operators $\bar\Omega^\dagger\Omega$ are evaluated with the expansion~\eqref{eq:OmegaX}: $\Omega$ acting on $\ket0$ involves only $\mf N$, $\mf A_2$ and $\mf B_2$; $\bar\Omega^\dagger$ then acts on states with up to two gluons. Note that $\bar\Omega^\dagger\Omega=1$ for $U=1$, so that every amplitude vanishes in the absence of the target.

From here on $\mf A\equiv\mf A^{(1)}_2$ and $\mf A^\dagger\equiv\mf A^{(1)\dagger}_2$; the unitarity relation $\mf A^{(1)\dagger}_1(p^+,\v x)=-\mf A(p^+,-\v x)$ is used to eliminate $\mf A^{(1)}_1$ wherever possible, and $\mf A^{(3)}\equiv\mf A^{(3)}_2$.

\subsection{Leading order}\label{sec:LO}

At order $g$ only $\mf A$ contributes:
\begin{align}\label{eq:ALO}
\mathcal A^{a}_{{\rm LO}\,i}(k)&=\int_{w}e^{-i\v k\cdot\v w}\Big[U^{ab}(\v w)\,\mf A^b_i(k^+,-\v w)-\bar{\mf A}^a_i(k^+,-\v w)\Big]
\notag\\&=\frac{ig}{\sqrt2\,\pi\sqrt{k^+}}\int_{w,x}e^{-i\v k\cdot\v w}\,\Kv^i(\v x-\v w)\big[U(\v w)-U(\v x)\big]^{ab}\rho^b(\v x)\,.
\end{align}
The first term is the emission of the gluon before the collision, followed by the scattering of the gluon; the second is the emission after the collision by the rotated charge. The LO spectrum is
\begin{align}\label{eq:LO}
\frac{d^3N_{\rm LO}}{d^3k}=\ &\frac{1}{(2\pi)^3}\,\frac{g^2}{2\pi^2}\,\frac{1}{k^+}\int_{x,x',w,w'}e^{-i\v k\cdot(\v w-\v w')}\,
\Kv(\v x-\v w)\cdot\Kv(\v x'-\v w')
\notag\\&\times
\Big\langle\big[U(\v x')-U(\v w')\big]^{ab'}\rho^{b'}(\v x')\,\big[U(\v x)-U(\v w)\big]^{ab}\rho^{b}(\v x)\Big\rangle\,.
\end{align}
Note that $\frac{1}{(2\pi)^3}\frac{g^2}{2\pi^2}=\frac{\alpha_s}{4\pi^4}$.

\paragraph{Colour-singlet projectile.} For a projectile that is a colour singlet, or more generally whose charge correlator is colour diagonal,
\begin{align}\label{eq:singlet}
\big\langle\rho^a(\v x)\rho^b(\v y)\big\rangle=\frac{\delta^{ab}}{N_c^2-1}\,n(\v x,\v y)\,,\qquad n(\v x,\v y)=\big\langle\rho^f(\v x)\rho^f(\v y)\big\rangle\,,
\end{align}
and the LO spectrum becomes
\begin{align}\label{eq:LOsinglet}
\frac{d^3N_{\rm LO}}{d^3k}=\ &\frac{\alpha_s}{4\pi^4(N_c^2-1)}\,\frac{1}{k^+}\int_{x,y,w,w'}e^{-i\v k\cdot(\v w-\v w')}\,\Kv(\v x-\v w)\cdot\Kv(\v y-\v w')
\notag\\&\times\,n(\v x,\v y)\,
\mathcal O_{\rm LO}(\v x,\v y;\v w,\v w')\,,
\end{align}
with the target operator
\begin{align}\label{eq:OLO}
\mathcal O_{\rm LO}(\v x,\v y;\v w,\v w')={\rm Tr}\Big\{\big[U^\dagger(\v x)-U^\dagger(\v w)\big]\big[U(\v y)-U(\v w')\big]\Big\}\,.
\end{align}
Since $U$ is real, ${\rm Tr}\{A^\dagger B\}={\rm Tr}\{B^\dagger A\}$ for adjoint matrices, and for a symmetric $n$ the integrand is invariant under $(\v x,\v w)\leftrightarrow(\v y,\v w')$; in the colour-singlet form the sign of the exponent is therefore immaterial. Eq.~\eqref{eq:LOsinglet} is the formula of Kovchegov and Tuchin~\cite{Kovchegov:2001sc}: for a single quark, $n\to C_F$ times a $\delta$-function, it reproduces their normalisation $\alpha_sC_F/(4\pi^4)$ exactly. It is also the operator that appears in the LL analysis of Kovner and Lublinsky~\cite{Kovner:2006wr,Kovner:2006ge}.

\paragraph{Momentum space and $k_T$ factorisation.} With $f(\v k)=\int_xe^{i\v k\cdot\v x}f(\v x)$ and Eq.~\eqref{eq:FT_K} of Appendix~\ref{app:fourier}, the LO amplitude takes the form
\begin{align}\label{eq:ALOmom}
\mathcal A^{a}_{{\rm LO}\,i}(k)=\frac{\sqrt2\,g}{\sqrt{k^+}}\int\frac{d^2q}{(2\pi)^2}\,\Big[\frac{k^i}{\v k^2}-\frac{q^i}{\v q^2}\Big]\,U^{ab}(\v q-\v k)\,\rho^b(-\v q)\,,\qquad
U^{ab}(\v l)=\int_xe^{i\v l\cdot\v x}U^{ab}(\v x)\,,
\end{align}
so that the gluon is produced with the Lipatov-type vertex $\v k/\v k^2-\v q/\v q^2$, where $\v q$ is the momentum supplied by the projectile and $\v k-\v q$ the momentum transferred by the target. The term $U^{ab}(\v l)\propto\delta^{ab}\delta^{(2)}(\v l)$ of a trivial target does not contribute, since the vertex vanishes at $\v q=\v k$. For a dilute target, $U^{ab}(\v x)=\delta^{ab}+f^{acb}\alpha^c(\v x)+O(\alpha^2)$, the amplitude is linear in $\alpha$, and with $\big|\frac{\v k}{\v k^2}-\frac{\v q}{\v q^2}\big|^2=\frac{(\v k-\v q)^2}{\v k^2\v q^2}$ the spectrum reduces to the $k_T$-factorised form~\cite{Gribov:1983ivg,Kovchegov:2001sc,Kharzeev:2003wz,Blaizot:2004wu}. For a projectile and a target that are homogeneous in the transverse plane, with $\langle\rho^{b'}(\v q')\rho^b(-\v q)\rangle=\delta^{bb'}(2\pi)^2\delta^{(2)}(\v q-\v q')P(\v q)$ and $\langle\alpha^{c}(\v l)\alpha^{d}(-\v l')\rangle=\delta^{cd}(2\pi)^2\delta^{(2)}(\v l-\v l')T(\v l)$,
\begin{align}\label{eq:kTfact}
\frac{d^3N_{\rm LO}}{d^3k}\simeq\frac{\alpha_s N_c(N_c^2-1)}{\pi^2}\,\frac{S_\perp}{k^+}\int\frac{d^2q}{(2\pi)^2}\,\frac{(\v k-\v q)^2}{\v k^2\,\v q^2}\,P(\v q)\,T(\v k-\v q)\,,
\end{align}
where $S_\perp$ is the transverse area.

\subsection{Next-to-leading order: amplitudes}\label{sec:NLOamplitudes}

At order $g^3$ the one-gluon amplitude receives five types of contributions. Collecting the terms of $\bar\Omega^\dagger\Omega$ that connect $\ket0$ with $\bra{g}$ at this order we find
\begin{align}\label{eq:A3amp}
\mathcal A^a_{(3)\,i}(k)=\int_w e^{-i\v k\cdot\v w}\Big\{\mathcal A_{N}+\mathcal A_{A^{(3)}}+\mathcal A_{B_3}+\mathcal A_{AB_2}+\mathcal A_{CB_2}\Big\}^a_i\,,
\end{align}
with
\begin{align}
\mathcal A_{N}&=-\bar{\mf A}^a_i(k^+,-\v w)\,\mf N_2+U^{ab}(\v w)\,\bar{\mf N}_2\,\mf A^b_i(k^+,-\v w)-\big[\bar{\mf N}_2,\bar{\mf A}^a_i(k^+,-\v w)\big]\,,
\notag\\
\mathcal A_{A^{(3)}}&=U^{ab}(\v w)\,\mf A^{(3)b}_i(k^+,-\v w)-\bar{\mf A}^{(3)a}_i(k^+,-\v w)\,,
\notag\\
\mathcal A_{B_3}&=\int\frac{dp^+}{2\pi}\int_x\bar{\mf B}^{\dagger ba}_{3ji}(k^+,\v w;p^+,-\v x)\Big[U^{bc}(\v x)\,\mf A^c_j(p^+,-\v x)-\bar{\mf A}^b_j(p^+,-\v x)\Big]\,,
\notag\\
\mathcal A_{AB_2}&=2\int\frac{dp^+}{2\pi}\int_x\Big[U^{ad}(\v w)U^{bc}(\v x)\,\bar{\mf A}^{\dagger b}_j(p^+,-\v x)\,\mf B^{cd}_{2ji}(k^+,-\v w;p^+,-\v x)
\notag\\&\hspace{2.2cm}
-\bar{\mf A}^{\dagger b}_j(p^+,-\v x)\,\bar{\mf B}^{ba}_{2ji}(k^+,-\v w;p^+,-\v x)\Big]\,,
\notag\\
\mathcal A_{CB_2}&=2\int\frac{dp^+}{2\pi}\int_{x,z}\Big[\bar{\mf C}^{\dagger dca}_{1jki}(k^+,\v w-\v z;p^+,\v z-\v x)\,U^{db}(\v x)U^{ce}(\v z)\,\mf B^{be}_{2jk}(k^+-p^+,-\v z;p^+,-\v x)
\notag\\&\hspace{2.2cm}
-\bar{\mf C}^{\dagger bca}_{1jki}(k^+,\v w-\v z;p^+,\v z-\v x)\,\bar{\mf B}^{bc}_{2jk}(k^+-p^+,-\v z;p^+,-\v x)\Big]\,.
\end{align}
The term $\mathcal A_N$ contains the normalisation: the term $\bar{\mf A}^{(3)}_1$ of $\bar\Omega^\dagger$ contributes, through the unitarity relation~\eqref{eq:A31unit}, the commutator $-[\bar{\mf N}_2,\bar{\mf A}]$. $\mathcal A_{A^{(3)}}$ is the $O(g^3)$ one-gluon amplitude of the wave function, before and after the collision. $\mathcal A_{B_3}$ describes the scattering of an already emitted gluon of momentum $p^+$ into the observed gluon by its interaction with the valence charges after the collision ($\bar{\mf B}_3^\dagger$); $\mathcal A_{AB_2}$, the absorption, after the collision, of one of the two gluons of the two-gluon component of the wave function; $\mathcal A_{CB_2}$, the recombination of the two gluons of that component into the observed one. In each term the two contributions correspond to an interaction of the soft gluons with the target before and after the respective valence interaction; their difference vanishes for $U=1$.

The terms of $\mathcal A_N$ are not all of the same order in $\rho$. With the ordering of the three-$\rho$ part of $\mf A^{(3)}$ derived in Section~\ref{sec:coefficients}, $[\mf A^{(3)}]_{\rho\rho\rho}=\mf A\,\mf N_2$, the commutator in $\mathcal A_N$ cancels against the corresponding part of $\mathcal A_{A^{(3)}}$, and the four-$\rho$ (normalisation) terms of the amplitude add up to
\begin{align}\label{eq:Nterms}
-\bar{\mf A}\,\mf N_2+U\,\bar{\mf N}_2\,\mf A+U\,\mf A\,\mf N_2-\bar{\mf N}_2\,\bar{\mf A}\,,
\end{align}
which vanishes for $U=1$, as it must. With the opposite ordering, $\mf N_2\mf A$, a spurious two-$\rho$ term $\propto\log(\vee/\Lambda)$ would remain.

The two-gluon amplitude at order $g^2$ is
\begin{align}\label{eq:A2gamp}
\mathcal A^{ab}_{ij}(k;q^+,\v v)=\frac{2}{\sqrt{2\pi}}\int_w e^{-i\v k\cdot\v w}\Big\{\mathcal A_B+\mathcal A^{(1)}_{AA}+\mathcal A^{(2)}_{AA}-\mathcal A_{CA}\Big\}^{ab}_{ij}\,,
\end{align}
with
\begin{align}
\mathcal A_B&=U^{ac}(\v w)\,\mf B^{bc}_{2ji}(k^+,-\v w;q^+,-\v v)-U^{\dagger bc}(\v v)\,\bar{\mf B}^{ca}_{2ji}(k^+,-\v w;q^+,-\v v)\,,
\notag\\
\mathcal A^{(1)}_{AA}&=\tfrac12\Big[U^{\dagger bc}(\v v)\,\bar{\mf A}^c_j(q^+,-\v v)\,\bar{\mf A}^a_i(k^+,-\v w)-U^{ac}(\v w)\,\bar{\mf A}^d_j(q^+,-\v v)\,U^{\dagger bd}(\v v)\,\mf A^c_i(k^+,-\v w)\Big]\,,
\notag\\
\mathcal A^{(2)}_{AA}&=\tfrac12\Big[U^{\dagger bc}(\v v)\,\bar{\mf A}^a_i(k^+,-\v w)\,\bar{\mf A}^c_j(q^+,-\v v)-\bar{\mf A}^a_i(k^+,-\v w)\,\mf A^b_j(q^+,-\v v)\Big]\,,
\notag\\
\mathcal A_{CA}&=\int_z\Big[\bar{\mf C}^{dac}_{1jik}(q^++k^+,\v z-\v w;q^+,\v w-\v v)\,U^{\dagger bd}(\v v)\,U^{ce}(\v z)\,\mf A^e_k(q^++k^+,-\v z)
\notag\\&\hspace{1.2cm}
-U^{\dagger bd}(\v v)\,\bar{\mf C}^{dac}_{1jik}(q^++k^+,\v z-\v w;q^+,\v w-\v v)\,\bar{\mf A}^c_k(q^++k^+,-\v z)\Big]\,.
\end{align}
$\mathcal A_B$ is the two-gluon component of the wave function scattered on the target; $\mathcal A^{(1)}_{AA}$ and $\mathcal A^{(2)}_{AA}$ arise from the product of two single emissions, one before and one after the collision, in the two operator orderings generated by $\bar\Omega^\dagger$; and $\mathcal A_{CA}$ is the splitting, after the collision, of a gluon of momentum $q^++k^+$ emitted before the collision into the observed gluon and a second gluon of momentum $q^+$.

\subsection{Next-to-leading order: Groups and rows}\label{sec:groups}

The $O(g^4)$ cross section is
\begin{align}\label{eq:NLOdecomp}
\frac{d^3N}{d^3k}=\frac{d^3N_{\rm LO}}{d^3k}+\frac{d^3N^{\rm(I)}}{d^3k}+\frac{d^3N^{\rm(II)}}{d^3k}+\frac{d^3N^{\rm(III)}}{d^3k}+O(g^6)\,,
\end{align}
where
\begin{align}\label{eq:Groups}
\frac{d^3N^{\rm(I)}}{d^3k}&=\frac{1}{(2\pi)^3}\Big\langle\mathcal A^{a\dagger}_{{\rm LO}\,i}(k)\,\mathcal A^a_{(3)\,i}(k)\Big\rangle+{\rm c.c.}\,,\notag\\
\frac{d^3N^{\rm(II)}}{d^3k}&=\frac{1}{(2\pi)^3}\int dq^+\!\int_v\Big\langle\big|\mathcal A_B\big|^2+\big|\mathcal A^{(1)}_{AA}\big|^2+\big|\mathcal A^{(2)}_{AA}\big|^2+\big|\mathcal A_{CA}\big|^2\Big\rangle\,,\notag\\
\frac{d^3N^{\rm(III)}}{d^3k}&=\frac{1}{(2\pi)^3}\int dq^+\!\int_v\Big\langle\mathcal A^{(1)\dagger}_{AA}\mathcal A_B+\mathcal A^{(2)\dagger}_{AA}\mathcal A_B-\mathcal A^{\dagger}_{CA}\mathcal A_B
\notag\\&\hspace{3.3cm}+\mathcal A^{(2)\dagger}_{AA}\mathcal A^{(1)}_{AA}-\mathcal A^{\dagger}_{CA}\mathcal A^{(1)}_{AA}-\mathcal A^{\dagger}_{CA}\mathcal A^{(2)}_{AA}\Big\rangle+{\rm c.c.}
\end{align}
Here $|\mathcal A|^2\equiv\mathcal A^{ab\dagger}_{ij}\mathcal A^{ab}_{ij}$, the sign of $\mathcal A_{CA}$ is that of Eq.~\eqref{eq:A2gamp}, and the momentum of the integrated gluon is renamed $q^+\to p^+$ below. Group~I is the interference of the LO amplitude with the $O(g^3)$ one-gluon amplitude; Group~II consists of the squares of the four components of the two-gluon amplitude; Group~III of their interferences. Groups~I and~III are written with their complex conjugates; Group~II is Hermitian by itself. We call the individual products of coefficients \emph{rows}. They are listed in Table~\ref{tab:rows}, together with the range of the longitudinal integration, the number of charges and the parts of the final result to which they contribute.

\begin{table}[t]
\centering
\renewcommand{\arraystretch}{1.2}
\footnotesize
\begin{tabular}{@{}l p{4.0cm} p{2.9cm} l p{3.4cm}@{}}
\toprule
Row & product (Eqs.~\eqref{eq:A3amp},~\eqref{eq:A2gamp}) & range of $p^+$ & $\#\rho$ & contributes to\\
\midrule
I.I & LO$^\dagger\times$ normalisation terms~\eqref{eq:Nterms} & --- & 4 & 4$\rho$ ($\log\frac{\vee}{\Lambda}$)\\
I.II & LO$^\dagger\times[U\mf A^{(3)}-\bar{\mf A}^{(3)}]$ & --- & 2--4 & 4$\rho$, BFKL, JIMWLK, 3$\rho$, finite\\
I.III & LO$^\dagger\times\bar{\mf B}_3^\dagger\mf A$ & $\Lambda<p^+<\vee$, $|p^+-k^+|>\Lambda$ & 3, 4 & BFKL, 3$\rho$, finite\\
I.IV & LO$^\dagger\times\bar{\mf A}^\dagger\mf B_2$ & $\Lambda<p^+<\vee-k^+$ $[\vee]$ & 3, 4 & 4$\rho$, BFKL, 3$\rho$; finite part $=-$III.I\\
I.V & LO$^\dagger\times\bar{\mf C}_1^\dagger\mf B_2$ & $\Lambda<p^+<k^+-\Lambda$ & 2, 3 & JIMWLK, 3$\rho$, finite\\
\midrule
II.I & $|\mathcal A_B|^2$ & $\Lambda<p^+<\vee-k^+$ $[\vee]$ & 2--4 & 4$\rho$, BFKL, JIMWLK, 3$\rho$, finite\\
II.II & $|\mathcal A^{(1)}_{AA}|^2$ & $\Lambda<p^+<\vee$ & 4 & 4$\rho$\\
II.III & $|\mathcal A^{(2)}_{AA}|^2$ & $\Lambda<p^+<\vee$ & 4 & 4$\rho$\\
II.IV & $|\mathcal A_{CA}|^2$ & $\Lambda<p^+<\vee-k^+$ & 2 & BFKL, JIMWLK, DGLAP, RC, finite\\
\midrule
III.I & $\mathcal A^{(1)\dagger}_{AA}\mathcal A_B$ & $\Lambda<p^+<\vee-k^+$ $[\vee]$ & 3, 4 & 4$\rho$, BFKL, 3$\rho$; finite part $=-$I.IV\\
III.II & $\mathcal A^{(2)\dagger}_{AA}\mathcal A_B$ & $\Lambda<p^+<\vee-k^+$ $[\vee]$ & 3, 4 & 4$\rho$, BFKL, 3$\rho$, finite\\
III.III & $-\mathcal A^{\dagger}_{CA}\mathcal A_B$ & $\Lambda<p^+<\vee-k^+$ & 2, 3 & BFKL, JIMWLK, 3$\rho$, finite\\
III.IV & $\mathcal A^{(2)\dagger}_{AA}\mathcal A^{(1)}_{AA}$ & $\Lambda<p^+<\vee$ & 4 & 4$\rho$\\
III.V & $-\mathcal A^{\dagger}_{CA}\mathcal A^{(1)}_{AA}$ & $\Lambda<p^+<\vee-k^+$ & 3 & BFKL, 3$\rho$, finite\\
III.VI & $-\mathcal A^{\dagger}_{CA}\mathcal A^{(2)}_{AA}$ & $\Lambda<p^+<\vee-k^+$ & 3 & BFKL, 3$\rho$, finite\\
\bottomrule
\end{tabular}
\caption{The rows of the $O(g^4)$ cross section. ``LO$^\dagger$'' is the conjugate LO amplitude, Eq.~\eqref{eq:ALO}; $p^+$ is the longitudinal momentum of the second (integrated or virtual) gluon. Where two ranges apply, the first refers to the part linear in $\rho$ of $\mf B_2$ (its parent gluon carries $p^++k^+$) and the bracket $[\vee]$ to its $\{\rho,\rho\}$ part, $\Lambda<p^+<\vee$. ``$\#\rho$'' is the number of charges before reordering; ``4$\rho$'' and ``3$\rho$'' denote the terms quartic and cubic in $\rho$ multiplying a rapidity logarithm, which cancel (Section~\ref{sec:rapidity}). Row I.II contains in addition the ultraviolet $\beta_0$ terms of the vertex correction, which cancel against Row I.V (Section~\ref{sec:dglap_rc}).}
\label{tab:rows}
\end{table}

The longitudinal ranges in Table~\ref{tab:rows} follow from the kinematics of the coefficients. All soft gluons satisfy $\Lambda<p^+<\vee$. Whenever a soft gluon of momentum $p^++k^+$ appears as an intermediate state (the parent of $\mf C_1$, the part of $\mf B_2$ linear in $\rho$), the condition $p^++k^+<\vee$ applies. The $\{\rho,\rho\}$ part of $\mf B_2$ describes two independent emissions and has support on the full square $\Lambda<k^+,p^+<\vee$; this is required for the cancellation of the four-$\rho$ terms (Section~\ref{sec:fourrho}). In $\mf B_3$ the momentum difference of the incoming and outgoing gluon is carried by the valence charges and satisfies $|p^+-k^+|>\Lambda$.

The evaluation of each row consists of the colour algebra of the products of coefficients, with the charges kept in their order, the longitudinal integration over $p^+$, and the transverse integrations that can be done in closed form. Three kinds of terms emerge:
\begin{enumerate}
\item terms proportional to $\log(\vee/k^+)$ or $\log(k^+/\Lambda)$, quadratic, cubic or quartic in $\rho$, which arise from the regions $p^+\gg k^+$ and $p^+\ll k^+$ of the longitudinal integrals. They are discussed in Section~\ref{sec:rapidity};
\item terms with a collinear singularity, from the region in which the integrated gluon of Row~II.IV is collinear with the observed one, and terms proportional to $\beta_0$ (Section~\ref{sec:dglap_rc});
\item the remaining finite terms (Section~\ref{sec:finite}).
\end{enumerate}
Two rows require a remark. The non-logarithmic parts of Rows~I.IV and III.I cancel identically. Rows~II.II, II.III and III.IV are products of four coefficients $\mf A$, each linear in $\rho$; they contain only four-$\rho$ terms.

% ======================================================================
\section{Rapidity logarithms: BFKL evolution of the projectile and JIMWLK evolution of the target}\label{sec:rapidity}
% ======================================================================

In this section we collect all terms of the $O(g^4)$ cross section that are enhanced by a logarithm of the longitudinal cut-offs and show that they are the leading-logarithmic evolution of the LO spectrum: the BFKL evolution of the projectile charge correlator over the rapidity interval between the projectile and the produced gluon, and the JIMWLK evolution of the target operator over the interval between the produced gluon and the target.

\subsection{Origin of the logarithms}\label{sec:loglong}

Rapidity logarithms arise from longitudinal integrals of the form $\int dp^+/p^+$ over the momentum of the second gluon, real (integrated over in Groups~II and~III) or virtual (inside the coefficients $\mf N_2$, $\mf A^{(3)}$, $\mf B_3$ and in the interference terms of Group~I). In every row the integrand behaves as $1/p^+$ in two regions: for $k^+\ll p^+<\vee$, where the second gluon is faster than the observed one and belongs to the projectile, and for $\Lambda<p^+\ll k^+$, where it is slower and is resolved by the observed gluon as part of the target. The first region produces $\ln(\vee/k^+)$, the second $\ln(k^+/\Lambda)$. Terms that do not depend on $k^+$ at all, in particular the normalisation $\mf N_2$ and the products of four single-emission amplitudes, carry the full logarithm
\begin{align}\label{eq:logsplit}
\ln\frac{\vee}{\Lambda}=\ln\frac{\vee}{k^+}+\ln\frac{k^+}{\Lambda}\,,
\end{align}
which we split at the longitudinal momentum $k^+$ of the produced gluon. With this split every logarithmic term is assigned either to the projectile or to the target side. Terms proportional to $\ln(1-k^+/\vee)$, which arise from the upper limit $\vee-k^+$ of the integrals with a soft parent gluon, vanish for $\vee\to\infty$ and are dropped.

By the order of the charges, the logarithmic terms are quartic, cubic or quadratic in $\rho$ (Table~\ref{tab:rows}). We discuss them in this order, because the reordering of the non-commuting charges feeds terms with more charges into terms with fewer.

\subsection{Terms quartic in \texorpdfstring{$\rho$}{rho}}\label{sec:fourrho}

Terms with four charges arise in Rows I.I, I.II (through the three-$\rho$ part $\mf A\mf N_2$ of $\mf A^{(3)}$), I.IV, II.I, II.II, II.III, III.I, III.II and III.IV, from the products of four single-emission coefficients $\mf A$, of $\mf A$ with $\mf N_2$, and of $\mf A$ with the $\{\rho,\rho\}$ part of $\mf B_2$. In all of them the longitudinal integral runs over the full interval $\Lambda<p^+<\vee$ and gives $\ln(\vee/\Lambda)$, split according to Eq.~\eqref{eq:logsplit}.

The totally symmetrised four-$\rho$ part of their sum vanishes identically, i.e.\ the four-$\rho$ terms cancel for commuting charges. This cancellation has a simple origin. For commuting charges the part of $\Omega$ that is linear in $\rho$ exponentiates: to the order considered, $\Omega$ acting on the vacuum is the coherent state
\begin{align}\label{eq:coherent}
\Omega_{\rm coh}\ket0=\exp\Big\{\int\frac{d^3p}{(2\pi)^3}\Big[\mf A^b_j(p)\,a^{\dagger b}_j(p)-\mf A^{b\dagger}_j(p)\,a^b_j(p)\Big]\Big\}\ket0\,,
\end{align}
whose expansion to second order produces $\mf N_2$ and the $\{\rho,\rho\}$ part of $\mf B_2$, $\mf B_2\supset\frac12\mf A\mf A$. Since $\bar\Omega_{\rm coh}=S^\dagger\Omega_{\rm coh}S$ is again a displacement operator, $\bar\Omega^\dagger_{\rm coh}\Omega_{\rm coh}$ is, up to a phase, a displacement operator, and the expectation value of the number operator in the corresponding coherent state is the squared modulus of the displacement. The displacement is the LO amplitude~\eqref{eq:ALO}. Hence for commuting charges the spectrum is exactly the LO result, with no correction of order $g^4\rho^4$. Two features of the coefficients are necessary for this. First, the $\{\rho,\rho\}$ part of $\mf B_2$ must have support on the full square $\Lambda<k^+,p^+<\vee$, as follows from the kinematics of two independent emissions; with the support $k^++p^+<\vee$ the four-$\rho$ terms of the different rows would carry different logarithms and would not cancel. Second, the three-$\rho$ part of $\mf A^{(3)}$ must be ordered as $\mf A\mf N_2$ (Section~\ref{sec:coefficients}), so that the normalisation terms combine into Eq.~\eqref{eq:Nterms}. We have verified the cancellation by an explicit term-by-term evaluation, with the charges represented by $SU(2)$ and $SU(3)$ matrices and the order of all operators kept.

For non-commuting charges the four-$\rho$ terms do not vanish as operators: rewriting each product in terms of totally symmetrised products leaves remainders with commutators. For two and three operators
\begin{align}\label{eq:symmetrise}
\rho_1\rho_2&=\tfrac12\{\rho_1,\rho_2\}+\tfrac12[\rho_1,\rho_2]\,,\notag\\
\rho_1\rho_2\rho_3&=S(\rho_1\rho_2\rho_3)+\tfrac12\Big(S\big([\rho_1,\rho_2]\rho_3\big)+S\big([\rho_1,\rho_3]\rho_2\big)+S\big([\rho_2,\rho_3]\rho_1\big)\Big)\notag\\&\quad+\tfrac16\Big(\big[[\rho_1,\rho_2],\rho_3\big]-\big[[\rho_2,\rho_3],\rho_1\big]\Big)\,,
\end{align}
where $S$ denotes the totally symmetrised product, $S(XY)=\frac12\{X,Y\}$, and similarly for four operators. By Eq.~\eqref{eq:rho} each commutator reduces the number of charges by one. The four-$\rho$ terms therefore leave a three-$\rho$ remainder, multiplied by $\ln(\vee/k^+)$ and by $\ln(k^+/\Lambda)$, and a two-$\rho$ remainder, which is part of the BFKL and JIMWLK terms below.

\subsection{Terms cubic in \texorpdfstring{$\rho$}{rho}}\label{sec:threerho}

Logarithmic terms with three charges arise directly in Rows I.II--I.V, II.I, III.I--III.III, III.V and III.VI, from the parts of $\mf B_2$, $\mf B_3$ and $\mf A^{(3)}$ that are linear or quadratic in $\rho$ multiplied by single emissions, and in addition from the reordering of the four-$\rho$ terms. They are of the form $if^{abc}\times$(real kernel)$\times$(Wilson lines)$\times\rho\rho\rho$, they change sign together with the sign of the exponent $e^{-i\v k\cdot(\v w-\v w')}$, and they are of the type of the ``first saturation corrections'' in the projectile.

For each of the two logarithms, the totally symmetrised three-$\rho$ part of the sum of all these terms vanishes identically, for an arbitrary function of $\v w$ and $\v w'$ in place of the phase. We have organised the sum by transverse kernel and by the number of Wilson lines (two, three and four) and verified the cancellation in each class separately, analytically and with explicit $SU(2)$ and $SU(3)$ matrices for the charges on a transverse lattice, so that contact terms are included. As operators the three-$\rho$ terms leave, by Eq.~\eqref{eq:symmetrise}, two-$\rho$ and one-$\rho$ remainders. The two-$\rho$ remainders are part of the BFKL and JIMWLK terms. The one-$\rho$ remainders are contact terms $\propto f^{abc}\rho^c(\v x)\delta^{(2)}(\v x-\v y)$, which do not contribute in the colour-singlet form.

The non-logarithmic three-$\rho$ terms do not cancel. They belong to the finite terms (Section~\ref{sec:finite}).

\subsection{The BFKL term}\label{sec:BFKL}

Collecting all two-$\rho$ terms proportional to $\ln(\vee/k^+)$, the genuine ones of the rows of Table~\ref{tab:rows} (I.II--I.IV, II.I, II.IV, III.I--III.III, III.V, III.VI) and the remainders of the reordered three- and four-$\rho$ terms, we find
\begin{align}\label{eq:BFKLrho}
\ln\frac{\vee}{k^+}\,\frac{d^{3}N_{\rm BFKL}}{d^{3}k} & = \frac{1}{(2\pi)^{3}}\,\frac{g^{4}}{16\pi^{5}}\,\frac{1}{k^{+}}\,\log\Big[\frac{\vee}{k^{+}}\Big]\notag\\&\times\int_{x,x',z,w,w'}e^{-ik\cdot (w-w')}\Bigg[ \notag\\
& +\mathcal{K}(x-w)\cdot\mathcal{K}(x'-w')\,\frac{1}{(x-z)^{2}}\Big[\notag\\
& -N_{c}\,U^{ab}(w')\,U^{ac}(w)\,\rho^{b}(x')\,\rho^{c}(x)\notag\\
&\quad +N_{c}\,U^{ab}(w')\,U^{ac}(x)\,\rho^{b}(x')\,\rho^{c}(x)\notag\\
& +N_{c}\,U^{ab}(x')\,U^{ac}(w)\,\rho^{b}(x')\,\rho^{c}(x) \notag\\
&\quad -N_{c}\,U^{ab}(x')\,U^{ac}(x)\,\rho^{b}(x')\,\rho^{c}(x)\Big]\notag\\
& +\mathcal{K}(x-w)\cdot\mathcal{K}(x'-w')\,\frac{1}{(x'-z)^{2}}\Big[\notag\\
& -N_{c}\,U^{ab}(w')\,U^{ac}(w)\,\rho^{b}(x')\,\rho^{c}(x)\notag\\
&\quad +N_{c}\,U^{ab}(w')\,U^{ac}(x)\,\rho^{b}(x')\,\rho^{c}(x)\notag\\
& +N_{c}\,U^{ab}(x')\,U^{ac}(w)\,\rho^{b}(x')\,\rho^{c}(x) \notag\\
&\quad -N_{c}\,U^{ab}(x')\,U^{ac}(x)\,\rho^{b}(x')\,\rho^{c}(x)\Big]\notag\\
& +\mathcal{K}(x-w)\cdot\mathcal{K}(x'-w')\,\mathcal{K}(x-z)\cdot\mathcal{K}(x'-z)\Big[\notag\\
& +2\,f^{abc}\,f^{dbe}\,U^{ha}(w')\,U^{hd}(w)\,\rho^{c}(x')\,\rho^{e}(x)\notag\\
&\quad -2\,f^{abc}\,f^{dbe}\,U^{ha}(w')\,U^{hd}(x)\,\rho^{c}(x')\,\rho^{e}(x)\notag\\
& -2\,f^{abc}\,f^{dbe}\,U^{ha}(x')\,U^{hd}(w)\,\rho^{c}(x')\,\rho^{e}(x) \notag\\
&\quad +2\,f^{abc}\,f^{dbe}\,U^{ha}(x')\,U^{hd}(x)\,\rho^{c}(x')\,\rho^{e}(x)\Big]\notag\\
& +\mathcal{K}(w-z)\cdot\mathcal{K}(w'-z)\,\mathcal{K}(x-z)\cdot\mathcal{K}(x'-z)\Big[\notag\\
& +2\,N_{c}\,\rho^{a}(x)\,\rho^{a}(x')\notag\\
&\quad -\,f^{abc}\,f^{dce}\,U^{ha}(w')\,U^{hd}(w)\,\rho^{b}(x)\,\rho^{e}(x')\notag\\
& -\,f^{abc}\,f^{dce}\,U^{ha}(w')\,U^{hd}(w)\,\rho^{b}(x')\,\rho^{e}(x) \notag\\
& +\,f^{abc}\,f^{dce}\,U^{ha}(w')\,U^{hd}(z)\,\rho^{b}(x)\,\rho^{e}(x')\notag\\
&\quad +\,f^{abc}\,f^{dce}\,U^{ha}(w')\,U^{hd}(z)\,\rho^{b}(x')\,\rho^{e}(x)\notag\\
& +\,f^{abc}\,f^{dce}\,U^{ha}(z)\,U^{hd}(w)\,\rho^{b}(x)\,\rho^{e}(x') \notag\\
&\quad +\,f^{abc}\,f^{dce}\,U^{ha}(z)\,U^{hd}(w)\,\rho^{b}(x')\,\rho^{e}(x)\Big]\notag\\
& +\mathcal{K}(x'-w')\cdot\mathcal{K}(w-z)\,\frac{1}{(x-z)^{2}}\Big[\notag\\
& -N_{c}\,U^{ab}(w')\,U^{ac}(w)\,\rho^{b}(x')\,\rho^{c}(x) \notag\\
& +N_{c}\,U^{ab}(w')\,U^{ac}(z)\,\rho^{b}(x')\,\rho^{c}(x)\notag\\
&\quad +N_{c}\,U^{ab}(x')\,U^{ac}(w)\,\rho^{b}(x')\,\rho^{c}(x)\notag\\
& -N_{c}\,U^{ab}(x')\,U^{ac}(z)\,\rho^{b}(x')\,\rho^{c}(x)\Big]\notag\\
& +\mathcal{K}(x'-w')\cdot\mathcal{K}(w-z)\,\mathcal{K}(x-z)\cdot\mathcal{K}(x'-z)\Big[\notag\\
& +2\,f^{abc}\,f^{cde}\,U^{hd}(w')\,U^{ha}(w)\,\rho^{b}(x)\,\rho^{e}(x')\notag\\
&\quad -2\,f^{abc}\,f^{cde}\,U^{hd}(w')\,U^{ha}(z)\,\rho^{b}(x)\,\rho^{e}(x')\notag\\
& -2\,f^{abc}\,f^{cde}\,U^{hd}(x')\,U^{ha}(w)\,\rho^{b}(x)\,\rho^{e}(x') \notag\\
&\quad +2\,f^{abc}\,f^{cde}\,U^{hd}(x')\,U^{ha}(z)\,\rho^{b}(x)\,\rho^{e}(x')\Big]\notag\\
& +\mathcal{K}(x-w)\cdot\mathcal{K}(w'-z)\,\frac{1}{(x'-z)^{2}}\Big[\notag\\
& -N_{c}\,U^{ab}(w')\,U^{ac}(w)\,\rho^{c}(x)\,\rho^{b}(x')\notag\\
&\quad +N_{c}\,U^{ab}(w')\,U^{ac}(x)\,\rho^{c}(x)\,\rho^{b}(x')\notag\\
& +N_{c}\,U^{ab}(z)\,U^{ac}(w)\,\rho^{c}(x)\,\rho^{b}(x') \notag\\
&\quad -N_{c}\,U^{ab}(z)\,U^{ac}(x)\,\rho^{c}(x)\,\rho^{b}(x')\Big]\notag\\
& +\mathcal{K}(x-w)\cdot\mathcal{K}(w'-z)\,\mathcal{K}(x-z)\cdot\mathcal{K}(x'-z)\Big[\notag\\
& +2\,f^{abc}\,f^{cde}\,U^{ha}(w')\,U^{hd}(w)\,\rho^{b}(x')\,\rho^{e}(x)\notag\\&\quad -2\,f^{abc}\,f^{cde}\,U^{ha}(w')\,U^{hd}(x)\,\rho^{b}(x')\,\rho^{e}(x)\notag\\
& -2\,f^{abc}\,f^{cde}\,U^{ha}(z)\,U^{hd}(w)\,\rho^{b}(x')\,\rho^{e}(x) \notag\\
&\quad +2\,f^{abc}\,f^{cde}\,U^{ha}(z)\,U^{hd}(x)\,\rho^{b}(x')\,\rho^{e}(x)\Big]\Bigg]
\end{align}
where the average over the projectile is implied. With non-commuting charges the projectile evolution produces the symmetrised product $\frac12\{\rho^b(\v x'),\rho^c(\v x)\}$ in the last four blocks; the fixed order written there differs from it by contact terms $\propto f^{abc}\rho^c(\v x)\delta^{(2)}(\v x-\v x')$, which do not contribute in the colour-singlet form. The prefactor is $\frac{1}{(2\pi)^3}\frac{g^4}{16\pi^5}=\frac{\alpha_s^2}{8\pi^6}$.

For a colour-singlet projectile, Eq.~\eqref{eq:singlet}, Eq.~\eqref{eq:BFKLrho} becomes
\begin{align}\label{eq:BFKLsinglet}
\ln\frac{\vee}{k^+}\,\frac{d^{3}N_{\rm BFKL}}{d^{3}k} & = \frac{1}{(2\pi)^{3}}\,\frac{g^{4}}{16\pi^{5}}\,\frac{1}{k^{+}}\,\frac{1}{N^{2}_{c}-1}\,\log\Big[\frac{\vee}{k^{+}}\Big]\,\int_{x,y,z,w,w'}e^{-ik\cdot (w-w')}\notag\\&\times\langle\rho^{f}(x)\rho^{f}(y) \rangle\,\Bigg[ \notag\\ & \mathcal K(x-w)\cdot\mathcal K(y-w')\,K(x,y,z)\Big[2 N_{c}\mathcal O_{\rm LO}(x,y;w,w')\Big]\notag\\
&+\mathcal K(x-w)\cdot\mathcal K(y-w')\,K(x,x,z)\Big[-N_{c}\mathcal O_{\rm LO}(x,y;w,w')\Big]\notag\\
&+\mathcal K(x-w)\cdot\mathcal K(y-w')\,K(y,y,z)\Big[-N_{c}\mathcal O_{\rm LO}(x,y;w,w')\Big]\notag\\
&+K(w,w',z)\,K(x,y,z)\Big[2 N_{c}\mathcal O_{\rm LO}(z,z;w,w')\Big]\notag\\
&+\mathcal K(x-w')\cdot\mathcal K(w-z)\,K(x,y,z)\Big[2 N_{c}\mathcal O_{\rm LO}(z,x;w,w')\Big]\notag\\
&+\mathcal K(x-w')\cdot\mathcal K(w-z)\,K(y,y,z)\Big[-N_{c}\mathcal O_{\rm LO}(z,x;w,w')\Big]\notag\\
&+\mathcal K(x-w)\cdot\mathcal K(w'-z)\,K(x,y,z)\Big[2 N_{c}\mathcal O_{\rm LO}(x,z;w,w')\Big]\notag\\
&+\mathcal K(x-w)\cdot\mathcal K(w'-z)\,K(y,y,z)\Big[-N_{c}\mathcal O_{\rm LO}(x,z;w,w')\Big]\Bigg]
\end{align}
with $K(\v x,\v y,\v z)$ of Eq.~\eqref{eq:kernels} and $\mathcal O_{\rm LO}$ of Eq.~\eqref{eq:OLO}. As for LO, the sign of the exponent is immaterial in this form.

Eq.~\eqref{eq:BFKLsinglet} is the LO spectrum~\eqref{eq:LOsinglet} with the charge correlator replaced by its one-step BFKL evolution over the rapidity interval $\ln(\vee/k^+)$:
\begin{align}\label{eq:BFKLasLO}
\ln\frac{\vee}{k^+}\,\frac{d^3N_{\rm BFKL}}{d^3k}=\ln\frac{\vee}{k^+}\;\frac{d^3N_{\rm LO}}{d^3k}\bigg|_{n(\v x,\v y)\to\frac{d}{dY}n(\v x,\v y)}\,,
\end{align}
where the evolution of the projectile charge correlator in the dilute regime is~\cite{Kovner:2006wr,Kovner:2006ge}
\begin{align}\label{eq:BFKLrhorho}
\frac{d}{dY}\big\langle\rho^f(\v x)\rho^f(\v y)\big\rangle=-\frac{\alpha_sN_c}{2\pi^2}\int_z\Big[&K_D(\v x,\v y,\v z)\,\big\langle\rho^f(\v x)\rho^f(\v y)\big\rangle
\notag\\&-2\,\delta^{(2)}(\v x-\v y)\int_{z'}K(\v z,\v z',\v x)\,\big\langle\rho^f(\v z)\rho^f(\v z')\big\rangle
\notag\\&
+\big(2K(\v y,\v z,\v x)-K(\v z,\v z,\v x)\big)\big\langle\rho^f(\v y)\rho^f(\v z)\big\rangle
\notag\\&+\big(2K(\v x,\v z,\v y)-K(\v z,\v z,\v y)\big)\big\langle\rho^f(\v x)\rho^f(\v z)\big\rangle\Big]\,,
\end{align}
with the dipole kernel $K_D(\v x,\v y,\v z)=K(\v x,\v x,\v z)-2K(\v x,\v y,\v z)+K(\v y,\v y,\v z)$. The first term of Eq.~\eqref{eq:BFKLrhorho} gives the first three lines of Eq.~\eqref{eq:BFKLsinglet} (virtual terms, the produced gluon is emitted by the original charges); the second term, in which the charge correlator is concentrated at the position $\v z$ of the new gluon, gives the fourth line (the produced gluon is emitted by the new gluon); and the last two terms give the remaining four lines (the produced gluon is emitted by the new gluon in the amplitude and by an original charge in the conjugate amplitude, or vice versa). In terms of the colour-singlet correlator of the dipole picture, Eq.~\eqref{eq:BFKLrhorho} is the linearised BK (dipole BFKL) equation; the check for a single colour dipole is described in Section~\ref{sec:checks}.

\subsection{The JIMWLK term}\label{sec:JIMWLK}

The two-$\rho$ terms proportional to $\ln(k^+/\Lambda)$, the genuine ones of Rows I.II, I.V, II.I, II.IV and III.III and the remainders of the reordered three- and four-$\rho$ terms, add up to
\begin{align}\label{eq:JIMWLKrho}
\ln\frac{k^+}{\Lambda}\,\frac{d^{3}N_{\rm JIMWLK}}{d^{3}k} & = \frac{1}{(2\pi)^{3}}\,\frac{g^{4}}{16\pi^{5}}\,\frac{1}{k^{+}}\,\log{\left(\frac{k^{+}}{\Lambda}\right)}\,\int_{x,x',z,w,w'}e^{-ik\cdot(w-w')}\,\Bigg[
\notag \\ & \mathcal{K}(x-w)\cdot\mathcal{K}(x'-w')\,\frac{1}{(x-z)^{2}}\Big[\notag\\
&+2\,N_{c}\,U^{ab}(w')\,U^{ac}(x)\,\rho^{b}(x')\,\rho^{c}(x) -2\,N_{c}\,U^{ab}(x')\,U^{ac}(x)\,\rho^{b}(x')\,\rho^{c}(x)\notag\\
& +2\,f^{abc}\,f^{deh}\,U^{da}(z)\,U^{el}(w')\,U^{hb}(x)\,\rho^{l}(x')\,\rho^{c}(x) \notag\\
&\quad -2\,f^{abc}\,f^{deh}\,U^{da}(z)\,U^{el}(x')\,U^{hb}(x)\,\rho^{l}(x')\,\rho^{c}(x)\Big]\notag\\
& +\mathcal{K}(x-w)\cdot\mathcal{K}(x'-w')\,\frac{1}{(x'-z)^{2}}\Big[\notag\\
& +2\,N_{c}\,U^{ab}(x')\,U^{ac}(w)\,\rho^{b}(x')\,\rho^{c}(x)\notag\\
&\quad -2\,N_{c}\,U^{ab}(x')\,U^{ac}(x)\,\rho^{b}(x')\,\rho^{c}(x)\notag\\
& +2\,f^{abc}\,f^{deh}\,U^{da}(z)\,U^{el}(w)\,U^{hb}(x')\,\rho^{c}(x')\,\rho^{l}(x) \notag\\
&\quad -2\,f^{abc}\,f^{deh}\,U^{da}(z)\,U^{el}(x)\,U^{hb}(x')\,\rho^{c}(x')\,\rho^{l}(x)\Big]\notag\\
& +\mathcal{K}(x-w)\cdot\mathcal{K}(x'-w')\,\mathcal{K}(x-z)\cdot\mathcal{K}(x'-z)\Big[\notag\\
& +2\,N_{c}\,U^{ab}(x')\,U^{ac}(x)\,\rho^{b}(x')\,\rho^{c}(x)\notag\\
&\quad +2\,f^{abc}\,f^{ade}\,U^{hb}(x')\,U^{hd}(x)\,\rho^{c}(x')\,\rho^{e}(x)\notag\\
& -2\,f^{abc}\,f^{deh}\,U^{da}(z)\,U^{eb}(x)\,U^{hl}(x')\,\rho^{l}(x')\,\rho^{c}(x) \notag\\
&\quad -2\,f^{abc}\,f^{deh}\,U^{da}(z)\,U^{eb}(x')\,U^{hl}(x)\,\rho^{c}(x')\,\rho^{l}(x)\Big]\notag\\
& +\mathcal{K}(x-w)\cdot\mathcal{K}(x'-w')\,\frac{1}{(w-z)^{2}}\Big[\notag\\
& -2\,N_{c}\,U^{ab}(w')\,U^{ac}(w)\,\rho^{b}(x')\,\rho^{c}(x)\notag\\
&\quad +2\,N_{c}\,U^{ab}(x')\,U^{ac}(w)\,\rho^{b}(x')\,\rho^{c}(x)\notag\\
& -2\,f^{abc}\,f^{deh}\,U^{da}(z)\,U^{el}(w')\,U^{hb}(w)\,\rho^{l}(x')\,\rho^{c}(x) \notag\\
&\quad +2\,f^{abc}\,f^{deh}\,U^{da}(z)\,U^{el}(x')\,U^{hb}(w)\,\rho^{l}(x')\,\rho^{c}(x)\Big]\notag\\
& +\mathcal{K}(x-w)\cdot\mathcal{K}(x'-w')\,\frac{1}{(w'-z)^{2}}\Big[\notag\\
& -2\,N_{c}\,U^{ab}(w')\,U^{ac}(w)\,\rho^{b}(x')\,\rho^{c}(x)\notag\\
&\quad +2\,N_{c}\,U^{ab}(w')\,U^{ac}(x)\,\rho^{b}(x')\,\rho^{c}(x)\notag\\
& -2\,f^{abc}\,f^{deh}\,U^{da}(z)\,U^{el}(w)\,U^{hb}(w')\,\rho^{c}(x')\,\rho^{l}(x) \notag\\
&\quad +2\,f^{abc}\,f^{deh}\,U^{da}(z)\,U^{el}(x)\,U^{hb}(w')\,\rho^{c}(x')\,\rho^{l}(x)\Big]\notag\\
& +\mathcal{K}(x-w)\cdot\mathcal{K}(x'-w')\,\mathcal{K}(w-z)\cdot\mathcal{K}(w'-z)\Big[\notag\\
& +2\,N_{c}\,U^{ab}(w')\,U^{ac}(w)\,\rho^{b}(x')\,\rho^{c}(x)\notag\\
&\quad +2\,f^{abc}\,f^{ade}\,U^{hb}(w')\,U^{hd}(w)\,\rho^{c}(x')\,\rho^{e}(x)\notag\\
& -2\,f^{abc}\,f^{deh}\,U^{da}(z)\,U^{eb}(w)\,U^{hl}(w')\,\rho^{l}(x')\,\rho^{c}(x) \notag\\
&\quad -2\,f^{abc}\,f^{deh}\,U^{da}(z)\,U^{eb}(w')\,U^{hl}(w)\,\rho^{c}(x')\,\rho^{l}(x)\Big]\notag\\
& +\mathcal{K}(x-w)\cdot\mathcal{K}(x'-w')\,\mathcal{K}(x-z)\cdot\mathcal{K}(w'-z)\Big[\notag\\
& -2\,N_{c}\,U^{ab}(w')\,U^{ac}(x)\,\rho^{b}(x')\,\rho^{c}(x)\notag\\
& -2\,f^{abc}\,f^{ade}\,U^{hb}(w')\,U^{hd}(x)\,\rho^{c}(x')\,\rho^{e}(x) \notag\\
&\quad +2\,f^{abc}\,f^{deh}\,U^{da}(z)\,U^{eb}(x)\,U^{hl}(w')\,\rho^{l}(x')\,\rho^{c}(x)\notag\\
& +2\,f^{abc}\,f^{deh}\,U^{da}(z)\,U^{eb}(w')\,U^{hl}(x)\,\rho^{c}(x')\,\rho^{l}(x)\Big]\notag\\
& +\mathcal{K}(x-w)\cdot\mathcal{K}(x'-w')\,\mathcal{K}(x'-z)\cdot\mathcal{K}(w-z)\Big[\notag\\
& -2\,N_{c}\,U^{ab}(x')\,U^{ac}(w)\,\rho^{b}(x')\,\rho^{c}(x)\notag\\&\quad -2\,f^{abc}\,f^{ade}\,U^{hb}(x')\,U^{hd}(w)\,\rho^{c}(x')\,\rho^{e}(x)\notag\\
& +2\,f^{abc}\,f^{deh}\,U^{da}(z)\,U^{eb}(w)\,U^{hl}(x')\,\rho^{l}(x')\,\rho^{c}(x) \notag\\
&\quad +2\,f^{abc}\,f^{deh}\,U^{da}(z)\,U^{eb}(x')\,U^{hl}(w)\,\rho^{c}(x')\,\rho^{l}(x)\Big]\Bigg]
\end{align}
For a colour-singlet projectile this becomes
\begin{align}\label{eq:JIMWLKsinglet}
    \ln\frac{k^+}{\Lambda}\,\frac{d^{3}N_{\rm JIMWLK}}{d^{3}k} & = \frac{1}{(2\pi)^{3}}\,\frac{g^{4}}{16\pi^{5}}\,\frac{1}{k^{+}}\,\frac{1}{N^{2}_{c}-1}\,\log{\left(\frac{k^{+}}{\Lambda}\right)}\,\int_{x,y,z,w,w'}e^{-ik\cdot(w-w')}\notag\\&\times\langle\rho^{f}(x)\rho^{f}(y) \rangle\,\Bigg[\notag \\ & \mathcal K(x-w)\cdot\mathcal K(y-w')\,K(w,w',z)\Big[4 N_{c}{\rm Tr}\{U^\dagger(w)U(w')\}\notag\\
&-2{\rm Tr}\{T^eU^\dagger(w)U(w')T^eU^\dagger(w')U(z)\}\notag\\
&\quad-2{\rm Tr}\{T^eU^\dagger(w)U(w')T^eU^\dagger(z)U(w)\}\Big]\notag\\
&+\mathcal K(x-w)\cdot\mathcal K(y-w')\,K(x,w',z)\Big[-4 N_{c}{\rm Tr}\{U^\dagger(x)U(w')\}\notag\\
&+2{\rm Tr}\{T^eU^\dagger(w')U(x)T^eU^\dagger(x)U(z)\}\notag\\
&\quad+2{\rm Tr}\{T^eU^\dagger(w')U(x)T^eU^\dagger(z)U(w')\}\Big]\notag\\
&+\mathcal K(x-w)\cdot\mathcal K(y-w')\,K(x,y,z)\Big[4 N_{c}{\rm Tr}\{U^\dagger(x)U(y)\}\notag\\
&-2{\rm Tr}\{T^eU^\dagger(x)U(y)T^eU^\dagger(y)U(z)\}\notag\\
&\quad-2{\rm Tr}\{T^eU^\dagger(x)U(y)T^eU^\dagger(z)U(x)\}\Big]\notag\\
&+\mathcal K(x-w)\cdot\mathcal K(y-w')\,K(y,w,z)\Big[-4 N_{c}{\rm Tr}\{U^\dagger(w)U(y)\}\notag\\
&+2{\rm Tr}\{T^eU^\dagger(w)U(y)T^eU^\dagger(y)U(z)\}\notag\\
&\quad+2{\rm Tr}\{T^eU^\dagger(w)U(y)T^eU^\dagger(z)U(w)\}\Big]\notag\\
&+\mathcal K(x-w)\cdot\mathcal K(y-w')\,K(w',w',z)\Big[\notag\\
&\quad-2 N_{c}{\rm Tr}\{U^\dagger(w)U(w')\}+2 N_{c}{\rm Tr}\{U^\dagger(x)U(w')\}\notag\\
&-{\rm Tr}\{T^eU^\dagger(w')U(x)T^eU^\dagger(w')U(z)\}\notag\\
&\quad-{\rm Tr}\{T^eU^\dagger(w')U(x)T^eU^\dagger(z)U(w')\}\notag\\
&+{\rm Tr}\{T^eU^\dagger(w)U(w')T^eU^\dagger(w')U(z)\}\notag\\
&\quad+{\rm Tr}\{T^eU^\dagger(w)U(w')T^eU^\dagger(z)U(w')\}\notag\\
&+{\rm Tr}\{T^eU^\dagger(w')U(x)\}\,{\rm Tr}\{T^eU^\dagger(w')U(z)\}\notag\\
&\quad+{\rm Tr}\{T^eU^\dagger(w)U(w')\}\,{\rm Tr}\{T^eU^\dagger(w')U(z)\}\Big]\notag\\
&+\mathcal K(x-w)\cdot\mathcal K(y-w')\,K(w,w,z)\Big[\notag\\
&\quad-2 N_{c}{\rm Tr}\{U^\dagger(w)U(w')\}+2 N_{c}{\rm Tr}\{U^\dagger(w)U(y)\}\notag\\
&+{\rm Tr}\{T^eU^\dagger(w)U(w')T^eU^\dagger(w)U(z)\}\notag\\
&\quad+{\rm Tr}\{T^eU^\dagger(w)U(w')T^eU^\dagger(z)U(w)\}\notag\\
&-{\rm Tr}\{T^eU^\dagger(w)U(y)T^eU^\dagger(w)U(z)\}\notag\\
&\quad-{\rm Tr}\{T^eU^\dagger(w)U(y)T^eU^\dagger(z)U(w)\}\notag\\
&-{\rm Tr}\{T^eU^\dagger(w)U(w')\}\,{\rm Tr}\{T^eU^\dagger(w)U(z)\}\notag\\
&\quad+{\rm Tr}\{T^eU^\dagger(w)U(y)\}\,{\rm Tr}\{T^eU^\dagger(w)U(z)\}\Big]\notag\\
&+\mathcal K(x-w)\cdot\mathcal K(y-w')\,K(x,x,z)\Big[\notag\\
&\quad2 N_{c}{\rm Tr}\{U^\dagger(x)U(w')\}-2 N_{c}{\rm Tr}\{U^\dagger(x)U(y)\}\notag\\
&-{\rm Tr}\{T^eU^\dagger(w')U(x)T^eU^\dagger(x)U(z)\}\notag\\
&\quad-{\rm Tr}\{T^eU^\dagger(w')U(x)T^eU^\dagger(z)U(x)\}\notag\\
&+{\rm Tr}\{T^eU^\dagger(x)U(y)T^eU^\dagger(x)U(z)\}\notag\\
&\quad+{\rm Tr}\{T^eU^\dagger(x)U(y)T^eU^\dagger(z)U(x)\}\notag\\
&-{\rm Tr}\{T^eU^\dagger(w')U(x)\}\,{\rm Tr}\{T^eU^\dagger(x)U(z)\}\notag\\
&\quad-{\rm Tr}\{T^eU^\dagger(x)U(y)\}\,{\rm Tr}\{T^eU^\dagger(x)U(z)\}\Big]\notag\\
&+\mathcal K(x-w)\cdot\mathcal K(y-w')\,K(y,y,z)\Big[\notag\\
&\quad2 N_{c}{\rm Tr}\{U^\dagger(w)U(y)\}-2 N_{c}{\rm Tr}\{U^\dagger(x)U(y)\}\notag\\
&-{\rm Tr}\{T^eU^\dagger(w)U(y)T^eU^\dagger(y)U(z)\}\notag\\
&\quad-{\rm Tr}\{T^eU^\dagger(w)U(y)T^eU^\dagger(z)U(y)\}\notag\\
&+{\rm Tr}\{T^eU^\dagger(x)U(y)T^eU^\dagger(y)U(z)\}\notag\\
&\quad+{\rm Tr}\{T^eU^\dagger(x)U(y)T^eU^\dagger(z)U(y)\}\notag\\
&-{\rm Tr}\{T^eU^\dagger(w)U(y)\}\,{\rm Tr}\{T^eU^\dagger(y)U(z)\}\notag\\
&\quad+{\rm Tr}\{T^eU^\dagger(x)U(y)\}\,{\rm Tr}\{T^eU^\dagger(y)U(z)\}\Big]\Bigg]
\end{align} 
where the products of traces with $T^e$ are the adjoint rewriting of the $ffUUU$ contractions of Eq.~\eqref{eq:JIMWLKrho}.

Eq.~\eqref{eq:JIMWLKrho} is the JIMWLK evolution of the target operator of the LO spectrum over the rapidity interval $\ln(k^+/\Lambda)$. With the JIMWLK Hamiltonian in the normalisation of~\cite{Lublinsky:2016meo},
\begin{align}\label{eq:HJIMWLK}
H_{\rm JIMWLK}=\frac{\alpha_s}{2\pi^2}\int_{x,y,z}K(\v x,\v y,\v z)\Big[J_L^e(\v x)J_L^e(\v y)+J_R^e(\v x)J_R^e(\v y)-2\,c^{ef}(\v z)\,J_L^e(\v x)J_R^f(\v y)\Big]\,,
\end{align}
where $J_L^e(\v x)$ and $J_R^e(\v x)$ act on a Wilson line at $\v x$ as $J_L^eU=T^eU$ and $J_R^eU=UT^e$, $c^{ef}(\v z)$ is defined by $U^\dagger(\v z)T^eU(\v z)=c^{ef}(\v z)T^f$, and $\frac{d}{dY}=-H_{\rm JIMWLK}$, we find
\begin{align}\label{eq:JIMWLKasLO}
\ln\frac{k^+}{\Lambda}\,\frac{d^3N_{\rm JIMWLK}}{d^3k}=-\ln\frac{k^+}{\Lambda}\;H_{\rm JIMWLK}\,\frac{d^3N_{\rm LO}}{d^3k}\,,
\end{align}
where $H_{\rm JIMWLK}$ acts on the Wilson lines of Eq.~\eqref{eq:LO} only. The prefactors match since $\frac{g^4}{16\pi^5}\big/\frac{g^2}{2\pi^2}=\frac{\alpha_s}{2\pi^2}$. Because $\mathcal O_{\rm LO}$ is bilinear in $U$, $H_{\rm JIMWLK}$ acting on it can be evaluated exactly, for an arbitrary (not necessarily colour-singlet or symmetric) ordered correlator $\langle\rho^{b'}(\v x')\rho^b(\v x)\rangle$; Eq.~\eqref{eq:JIMWLKrho} agrees with the result pointwise in all transverse coordinates, and Eq.~\eqref{eq:JIMWLKsinglet} with its colour-singlet projection (Section~\ref{sec:checks}). For a single adjoint dipole the action of $H_{\rm JIMWLK}$ reduces to the BK form, $-H_{\rm JIMWLK}{\rm Tr}[U^\dagger(\v a)U(\v b)]=\frac{\alpha_s}{\pi^2}\int_zK_D(\v a,\v b,\v z)\big[c^{ef}(\v z){\rm Tr}(U^\dagger(\v a)T^eU(\v b)T^f)-N_c{\rm Tr}(U^\dagger(\v a)U(\v b))\big]$.

\subsection{Why the logarithms are the evolution of the LO spectrum}\label{sec:proofsketch}

The identifications~\eqref{eq:BFKLasLO} and~\eqref{eq:JIMWLKasLO} are the expansion to first order of the leading-logarithmic factorisation of the single inclusive spectrum~\cite{Kovner:2006wr,Kovner:2006ge,Gelis:2008rw}: the spectrum of a gluon at rapidity $y$ is the LO expression evaluated with the projectile evolved from its own rapidity $Y_P$ down to $y$ and with the target evolved from $y$ to its own rapidity,
\begin{align}\label{eq:LLfact}
\frac{d^3N}{d^3k}\bigg|_{\rm LL}=\Big\langle\,\frac{d^3N_{\rm LO}}{d^3k}\,\Big\rangle_{P,\,Y_P-y;\;T,\,y-Y_T}\,.
\end{align}
The structure of the calculation shows why this holds, and we sketch the argument. Consider first the region $p^+\gg k^+$. Here the second gluon is much faster than the observed one, and its emission and absorption amplitudes factorise from those of the observed gluon: in all coefficients the dependence on $k^+$ disappears at leading power in $k^+/p^+$, and the observed gluon is emitted by the total colour charge of the valence charges and of the fast gluon, $\rho^a(\v x)\to\rho^a(\v x)+$ (charge density of the gluon at $\v z$). The sum over the rows is then the LO spectrum evaluated in the projectile state that contains one more soft gluon, minus the probability that it is not there (the virtual terms with $\mf N_2$ and the self-energy parts of $\mf A^{(3)}$), i.e.\ the LO spectrum with the projectile correlator evolved by one step. For a dilute projectile this step is linear in $\langle\rho\rho\rangle$, Eq.~\eqref{eq:BFKLrhorho}; terms with three and four charges, which would describe the non-linear (KLWMIJ) evolution of a dense projectile, must and do cancel at this order in the symmetrised sense (Sections~\ref{sec:fourrho},~\ref{sec:threerho}). In the region $p^+\ll k^+$ the second gluon is much slower than the observed one. It is emitted coherently by the valence charges and by the observed gluon, and at leading power its emission does not change the momentum of the observed gluon. The rows then combine into the change of the scattering amplitudes of the valence charges and of the observed gluon due to one soft gluon that scatters together with them, i.e.\ into one step of the JIMWLK evolution of the Wilson lines of the LO amplitude and of its conjugate, including the virtual terms. Both statements are verified by the explicit results~\eqref{eq:BFKLrho}--\eqref{eq:JIMWLKsinglet}.

Two features of the calculation enter this argument in an essential way. First, the logarithms are split at the rapidity of the observed gluon, Eq.~\eqref{eq:logsplit}; the four-$\rho$ terms and the normalisation, which carry $\ln(\vee/\Lambda)$, are distributed between the two evolutions accordingly. Second, the operator ordering of the charges matters: it is the commutator terms of the reordering that produce part of the two-$\rho$ structures of both evolutions.

\subsection{Summary of the logarithmic terms}\label{sec:logsummary}

The rapidity-enhanced part of the $O(g^4)$ cross section is
\begin{align}\label{eq:logsummary}
\frac{d^3N}{d^3k}\bigg|_{\ln}=\ln\frac{\vee}{k^+}\,\frac{d^3N_{\rm BFKL}}{d^3k}+\ln\frac{k^+}{\Lambda}\,\frac{d^3N_{\rm JIMWLK}}{d^3k}\,,
\end{align}
with the BFKL term~\eqref{eq:BFKLrho} and the JIMWLK term~\eqref{eq:JIMWLKrho}, and no term with three or four charges multiplies a rapidity logarithm. Individual rows contain also double logarithms, $\ln^2(k^+/\Lambda)$, from the region in which the second gluon is both soft and far from the observed one in the transverse plane. They come from the recoil of the observed gluon in the rows with $\mf B_3$ and $\mf C_1$ and cancel in the sum over rows (Section~\ref{sec:fin_soft}).
% Source of the ln^2 cancellation: limits_checked/Cancellation_of_flagged_terms.tex (numerical, SU(2), SU(3)), presented in sec_finite.

% ======================================================================
\section{Collinear (DGLAP) and running-coupling terms}\label{sec:dglap_rc}
% ======================================================================
% SLOT. Input: paper/notes/fix_poles_dglap.tex (problems P3, P4 of BRIEF_fix.md), to be written
% by the agent fixing P3/P4. Further inputs: final_checked/Final_terms_checked.tex secs. 3-4
% (N_DGLAP with phase e^{-ik.(y-y')/xi}, N_RC, LO+RC), group_II_checked/Group_II_Row_IV_checked.tex,
% limits_checked/Divergent_terms_checked.tex, notes/lit_dglap_rc.tex.
% Content expected here:
%  - origin: Group II Row IV, final-state g -> gg splitting of the LO gluon, xi = k^+/(p^+ + k^+);
%  - account of all 1/epsilon poles: UV of the vertex correction (11/3 in Row I.II) and the
%    -11/6 bracket of Row I.V (C_1^dagger B_2, "2b UV"), their exact cancellation; collinear pole
%    of the real splitting absorbed into the fragmentation function; character of the RC pole;
%  - DGLAP term with the subtraction inside the splitting function, [(1-xi)/xi]_+ and xi/(1-xi)_+,
%    the -1/k^+ x (BFKL structure) remnant, the MS-bar (CXY) form of the logarithm (P4);
%  - the running-coupling term as the beta_0 delta(1-xi)/2 part of P_gg (add and subtract),
%    LO + RC = g^2(mu) x LO, scale c_0^2/r^2; relation to KLSZ 2308.15545.
% Labels used elsewhere: sec:dglap_rc.

\pending{Placeholder.} This section will present the collinear term, which arises from Row~II.IV of Table~\ref{tab:rows} when the integrated gluon is collinear with the observed one, the account of the ultraviolet and collinear poles of the calculation, and the running-coupling term, defined as the $\frac{\beta_0}{2}\delta(1-\xi)$ part of the gluon--gluon splitting function. Its main statements and their comparison with the literature are anticipated in Section~\ref{sec:lit_dglap}.

% ======================================================================
\section{Finite terms}\label{sec:finite}
% ======================================================================
% SLOT. Inputs: paper/notes/fix_soft.tex (problems P1, P2, P5: Group I Row III recoil term and
% its cancellation partner, soft remainders of the plus prescription and their virtual partners,
% large-distance tails, explicit ln|.|, distributional form of the instantaneous Group III terms)
% and paper/notes/fix_singlet.tex (problem P6: finite terms averaged over a Gaussian/colour-singlet
% projectile, <rho^a(x) rho^b(y)> = delta^{ab} n(x,y)/(N_c^2-1), colour-singlet form with traces
% of adjoint Wilson lines; cubic terms odd in k, vanishing for a Gaussian projectile).
% Further inputs: final_checked/Final_terms_checked.tex secs. 5-8, limits_checked/Finite_terms_final.tex,
% limits_checked/Cancellation_of_flagged_terms.tex (ln^2(k^+/Lambda) cancels in the sum over rows).
% Conventions to keep: phase e^{-ik.(w-w')}, c.c. for Groups I and III, principal values,
% plus prescription [1/p^+]_+, Lambda -> 0 and vee -> infinity, Group I Row IV and Group III Row I
% cancel identically, the 67/9 - pi^2/3 term of Row I.II (eq:constantLO).
% Labels used elsewhere: sec:finite.

\pending{Placeholder.} This section will list the finite terms of Groups~I, II and III, row by row (Table~\ref{tab:rows}), in the limit $\Lambda\to0$, $\vee\to\infty$, both for an arbitrary ordered charge correlator and in colour-singlet form. It will include the cancellation of the double logarithms $\ln^2(k^+/\Lambda)$ between the rows with recoil of the observed gluon, the treatment of the soft remainders of the plus prescriptions and of their large-distance behaviour, and the explicit distributional form of the instantaneous terms of Group~III. The term of Row~I.II proportional to the LO structure, Eq.~\eqref{eq:constantLO}, is already given in Section~\ref{sec:lit_constant}.

% ======================================================================
\section{Limits, consistency checks and comparison with the literature}\label{sec:checks}
% ======================================================================
% Sources: paper/notes/README.md, notes/lit_dglap_rc.tex, notes/lit_finite_dilute.tex,
% inputs/reviews/review_2.txt (secs. 3-6, 11, 13), limits_checked/*.
% Items marked with \pending{...} depend on the open fixes P1-P6 (BRIEF_fix.md) and must be
% revised when sections dglap_rc and finite are filled in.

The calculation involves a large number of terms, with non-commuting charges and Wilson lines in a definite order. In this section we list the checks to which the result has been subjected and compare it with known results. Unless stated otherwise, the algebraic identities quoted below were verified both analytically and numerically: the expressions were evaluated on a transverse lattice with random $SU(2)$ and $SU(3)$ Wilson lines in the adjoint representation and with either a completely general ordered charge correlator $\langle\rho^b(\v u)\rho^c(\v v)\rangle$ or the colour-singlet one, Eq.~\eqref{eq:singlet}. Since the Wilson lines are random, an identity that holds on the lattice to round-off error cannot hold by accident.

\subsection{Limits}\label{sec:limits}

\paragraph{No target, pure-gauge target.} For $U=1$ the operator $\bar\Omega^\dagger\Omega$ is the identity and every amplitude vanishes. In the result this is manifest: every term of every row, and every term of the BFKL, JIMWLK, DGLAP and running-coupling contributions, contains a factor that vanishes at $U=1$, such as $[U(\v x')-U(\v w')]^{ab'}\rho^{b'}(\v x')$, its analogue with three Wilson lines, or a combination of $f^{abc}$ and Wilson lines that reduces to $N_c-N_c-N_c+N_c=0$. A constant $U$ is a global colour rotation, which cannot produce gluons; all terms vanish for constant $U$ as well, by $f^{abc}U^{bd}U^{ce}=U^{ah}f^{hde}$ and $U^TU=1$.

\paragraph{Counting of charges.} The BFKL and JIMWLK terms are of order $g^4\rho^2$. No term of order $g^4\rho^4$ survives (Section~\ref{sec:fourrho}), and the three-$\rho$ terms that multiply a rapidity logarithm cancel in the symmetrised sense (Section~\ref{sec:threerho}). The finite terms are of order $g^4\rho^2$ and $g^4\rho^3$.

\paragraph{Reality.} Groups I and III are written with their complex conjugates and are real. Group II is Hermitian by itself: its kernels are symmetric under $(\v x,\v w)\leftrightarrow(\v x',\v w')$ and so is its colour operator, which we verified numerically to $10^{-12}$.

\paragraph{Odd terms.} Terms of the form $if^{abc}\times$(real kernel)$\times$(real colour structure)$\times e^{-i\v k\cdot(\v w-\v w')}$ give, together with their complex conjugate, $2\sin[\v k\cdot(\v w-\v w')]\times$(real), which is odd under $\v k\to-\v k$. All three-$\rho$ terms are of this type; they vanish for a Gaussian (colour-singlet, McLerran--Venugopalan-type) average over the projectile, which has no odd moments. This is the behaviour of the ``first saturation correction'' found in~\cite{McLerran:2016snu,Kovchegov:2018jun}: odd in $\v k$ configuration by configuration and zero after the average for the single inclusive spectrum. \pending{To be restated with the colour-averaged finite terms of Section~\ref{sec:finite} (input P6).}

\paragraph{Scaling.} Under a dilatation of all transverse coordinates, $\v x\to\lambda\v x$, $\v k\to\v k/\lambda$, every finite term scales like the LO spectrum. The finite terms contain no renormalisation scale; the dependence on $\mu$ is confined to the collinear and running-coupling terms (Section~\ref{sec:dglap_rc}).

\paragraph{Dilute target.} For $U=1+O(\alpha)$ the LO spectrum reduces to the $k_T$-factorised form, Eq.~\eqref{eq:kTfact}, with the Lipatov-type vertex. In this limit the BFKL and JIMWLK terms reduce to the LO BFKL evolution of the projectile and of the target unintegrated gluon distributions, respectively. A term-by-term comparison of the finite terms with the one-loop gluon production vertex in $k_T$ factorisation~\cite{Fadin:1998py,DelDuca:1998cx,Bartels:2006hg} requires all finite terms in momentum space and has not been carried out; see Section~\ref{sec:lit_constant} for the comparison of the constant.

\subsection{Leading order}

The LO spectrum, Eq.~\eqref{eq:LOsinglet}, coincides with the result of Kovchegov and Tuchin~\cite{Kovchegov:2001sc}, including its normalisation: for a quark projectile, $n\to C_F$, the prefactor is $\alpha_sC_F/(4\pi^4)$. The target operator $\mathcal O_{\rm LO}$ is the one of Kovner, Lublinsky and Weigert~\cite{Kovner:2006wr,Kovner:2006ge}.

\subsection{Projectile evolution: the BFKL term}\label{sec:check_BFKL}

\paragraph{Internal consistency.} The ordered form~\eqref{eq:BFKLrho} is an exact regrouping of the term-by-term list of all two-$\rho$ terms proportional to $\ln(\vee/k^+)$, for an arbitrary ordered correlator (relative difference $\lesssim10^{-15}$ for $SU(2)$, $10^{-14}$ for $SU(3)$), and Eq.~\eqref{eq:BFKLsinglet} is its colour-singlet reduction to the same accuracy.

\paragraph{Exact one-gluon evolution of a colour-singlet projectile.} As an independent test we computed the change of the spectrum due to one additional gluon with $k^+<p^+<\vee$ in a colour-singlet projectile directly, without the expansion in $\rho$ used in the rest of the paper. Any set of colour sources $s$ at positions $\v x_s$, with charges $Q_s^e$, produces the observed gluon with the operator $\sum_{s,s'}Q^e_{s'}Q^e_sF(\v x_s,\v x_{s'})$, where
\begin{align}
F(\v a,\v b)=\int_{w,w'}e^{-i\v k\cdot(\v w-\v w')}\Kv(\v a-\v w)\cdot\Kv(\v b-\v w')\,{\rm Tr}\big[(U^\dagger(\v a)-U^\dagger(\v w))(U(\v b)-U(\v w'))\big]\,,
\end{align}
since the colour rotation of the sources cancels in the inclusive sum. For a quark--antiquark dipole at $\v x_1$, $\v x_2$ the three charges of the $q\bar qg$ colour singlet satisfy $Q_1^2=Q_2^2=C_F$, $Q_g^2=N_c$, $Q_1\cdot Q_2=1/(2N_c)$ and $Q_{1,2}\cdot Q_g=-N_c/2$, and the probability to emit the additional gluon at $\v z$ is $\frac{g^2}{4\pi^3}\ln\frac{\vee}{k^+}C_F\,K_D(\v x_1,\v x_2,\v z)$. Real minus virtual gives, exactly at finite $N_c$,
\begin{align}
\delta\Big[\frac{d^3N}{d^3k}\Big]&=\frac{\alpha_sN_c}{2\pi^2}\ln\frac{\vee}{k^+}\int_zK_D(\v x_1,\v x_2,\v z)\big[\mathcal D_{1z}+\mathcal D_{z2}-\mathcal D_{12}\big]\times(\text{LO prefactor})\,,\notag\\
\mathcal D_{ab}&=\frac{1}{2N_c}\big[F_{aa}+F_{bb}-F_{ab}-F_{ba}\big]\,,
\end{align}
the dipole BFKL (linearised BK) kernel acting on the projectile. Substituting $n=C_F[\delta_{11}+\delta_{22}-\delta_{12}-\delta_{21}]$ into Eq.~\eqref{eq:BFKLsinglet} reproduces this result pointwise in $\v z$. Since every symmetric $n(\v x,\v y)$ with vanishing row sums, i.e.\ every colour-singlet projectile, is a linear combination of such dipole correlators, and Eq.~\eqref{eq:BFKLsinglet} is linear in $n$, the BFKL term is the exact leading-logarithmic projectile evolution for every colour-singlet projectile. The calculation with fundamental charges is exact to all orders in the charges, which shows that the expansion in $\rho$ loses nothing for a dipole.

\paragraph{Literature.} This is the structure found by Kovner and Lublinsky for the single inclusive spectrum~\cite{Kovner:2006wr} and by Kovner, Lublinsky and Weigert~\cite{Kovner:2006ge}: the projectile correlator evolves with BFKL from the projectile rapidity down to the rapidity of the produced gluon. It is the first-order expansion of the leading-logarithmic factorisation theorem of Gelis, Lappi and Venugopalan~\cite{Gelis:2008rw}, and it satisfies the general structure of the NLO correction discussed by Gelis and Laidet~\cite{Gelis:2012ri}.

\subsection{Target evolution: the JIMWLK term}\label{sec:check_JIMWLK}

We computed $-H_{\rm JIMWLK}\,\mathcal O_{\rm LO}$ independently, with the Hamiltonian~\eqref{eq:HJIMWLK} in the normalisation of~\cite{Lublinsky:2016meo,Kovner:2013ona}, for a general ordered correlator. The term-by-term list of the two-$\rho$ terms proportional to $\ln(k^+/\Lambda)$, the grouped form~\eqref{eq:JIMWLKrho} and $-H_{\rm JIMWLK}$ acting on the LO spectrum agree to $10^{-14}$ for $SU(2)$ and $SU(3)$, and the colour-singlet form~\eqref{eq:JIMWLKsinglet} agrees with the colour-singlet projection to $10^{-15}$. The JIMWLK term is therefore the JIMWLK evolution of the LO spectrum, with the correct sign and normalisation, for an arbitrary projectile and not only for a colour-singlet one. Together with Section~\ref{sec:check_BFKL}, this confirms at fixed order the factorisation of the rapidity logarithms at the rapidity of the produced gluon~\cite{Kovner:2006wr,Kovner:2006ge,Gelis:2008rw}.

\subsection{Wave function}\label{sec:check_wf}

The light-cone wave function from which the coefficients of $\Omega$ are obtained was computed to order $g^3$, with the same longitudinal cut-offs ($\vee=\Lambda e^{\delta Y}$) and $d_\perp=2-\epsilon$, by Lublinsky and Mulian~\cite{Lublinsky:2016meo}. Our components agree with theirs as follows:
\begin{center}
\renewcommand{\arraystretch}{1.2}
\small
\begin{tabular}{@{}p{4.6cm}p{2.3cm}p{6.6cm}@{}}
\toprule
component & Ref.~\cite{Lublinsky:2016meo} & comparison\\
\midrule
LO soft gluon, normalisation $\mf N_2$ & eqs.~2.50--2.52 & identical\\
two gluons, $O(g^2)$ & eqs.~3.95--3.96 & identical\\
one gluon, $O(g^3)$, linear in $\rho$ & eqs.~3.97--3.98 & identical for $N_f=0$, incl.\ $\frac{11}{3}$, $\frac{67}{9}-\frac{\pi^2}{3}$, $\ln^2$ cut-off terms\\
one gluon, $O(g^3)$, $f^{abc}\{\rho^b,\rho^c\}$ & eqs.~3.100--3.101 & equal after the antisymmetrisation enforced by $f^{abc}\{\rho^b,\rho^c\}$\\
one gluon, $O(g^3)$, $\rho\rho\rho$ & eq.~3.102 & identical up to the operator ordering, see below\\
\bottomrule
\end{tabular}
\end{center}
In the two-$\rho$ component the two forms differ by terms symmetric under the exchange of the two charges, which do not contribute. The coordinate-space forms of the $\{\rho,\rho\}$ component, Eq.~\eqref{eq:A3tworho}, and of the coefficients go beyond the published material. The operator ordering of the three-$\rho$ component, $\rho^a(\v w)\rho^b(\v x)\rho^b(\v y)$, i.e.\ $[\mf A^{(3)}]_{\rho\rho\rho}=\mf A\,\mf N_2$, follows from third-order perturbation theory with operator-valued matrix elements; we verified it against a construction of $\Omega$ and an exact diagonalisation with finite-dimensional matrices representing the charges (agreement to $3\times10^{-16}$; the opposite ordering deviates at order one).

\subsection{The constant \texorpdfstring{$\frac{67}{9}-\frac{\pi^2}{3}$}{67/9-pi2/3}}\label{sec:lit_constant}

The term of Row I.II proportional to the LO structure is, with its complex conjugate,
\begin{align}\label{eq:constantLO}
\frac{\alpha_sN_c}{2\pi}\Big(\frac{67}{9}-\frac{\pi^2}{3}\Big)\frac{d^3N_{\rm LO}}{d^3k}\,.
\end{align}
It follows directly from the one-$\rho$ part of $\mf A^{(3)}$, Eq.~\eqref{eq:A3onerho}, and the normalisation of the LO amplitude. Before the cancellation of the ultraviolet $\beta_0$ logarithms against Row I.V, Row I.II with its complex conjugate is $\frac{\alpha_sN_c}{2\pi}\big[\frac{11}{3}\big(2\gamma_E+\frac12\ln\frac{X^2\mu^2_{\MS}}{4}+\frac12\ln\frac{X'^2\mu^2_{\MS}}{4}\big)+\frac{67}{9}-\frac{\pi^2}{3}\big]\times$LO, $X=x-w$, $X'=x'-w'$. This is, term by term at $N_f=0$, the bare one-gluon kernel $K_{JSJ}$ of the NLO JIMWLK Hamiltonian in the form of Kovner, Lublinsky, Skokov and Zhao~\cite{Kovner:2023vsy}, with the same sign, size and scale; since both are built from the wave function of~\cite{Lublinsky:2016meo}, this is a consistency check rather than an independent comparison. The inclusive quantities in which the same constant appears, the NLO BK kernel of Balitsky and Chirilli~\cite{Balitsky:2007feb}, the NLO BFKL kernel of Fadin and Lipatov~\cite{Fadin:1998py} and the NLO jet vertex of Bartels, Sabio Vera and Schwennsen~\cite{Bartels:2006hg}, have the coefficient $\frac{\alpha_sN_c}{4\pi}(\frac{67}{9}-\frac{\pi^2}{3})$, half of Eq.~\eqref{eq:constantLO}. These quantities are inclusive over the collinear gluon pair; as noted in~\cite{Kovner:2023vsy}, half of the constant sits there in the two-gluon terms. For the single inclusive gluon, whose collinear splitting is factorised into the fragmentation function (Section~\ref{sec:dglap_rc}), the bare-vertex value~\eqref{eq:constantLO} is the appropriate one, and the other half must reside in the collinear finite terms. \pending{Not verified: the momentum-weighted ($\int d\xi\,\xi$) inclusive combination of all collinear terms should give $-\frac{\alpha_sN_c}{4\pi}(\frac{67}{9}-\frac{\pi^2}{3})\times$LO; depends on the final form of the collinear and finite terms (P3, P4).} The amplitude-level one-loop Lipatov vertex~\cite{DelDuca:1998cx} contains no $\frac{67}{9}$ (it appears only after the real two-gluon production is added) and cannot be compared term by term with the light-cone wave-function constant.

\subsection{Collinear and running-coupling terms}\label{sec:lit_dglap}
% Depends on P3, P4 (input notes/fix_poles_dglap.tex). Revise after sec_dglap_rc is filled.

\pending{This subsection summarises the comparison of notes/lit\_dglap\_rc.tex and must be revised when Section~\ref{sec:dglap_rc} is written (P3, P4).} The collinear term comes from Row II.IV, in which the observed gluon is produced in the splitting $g\to gg$ of a gluon of momentum $k^+/\xi$ that was produced at LO; it is a final-state splitting. Compared with the forward hadron production results of Chirilli, Xiao and Yuan~\cite{Chirilli:2011km,Chirilli:2012jd} and of Altinoluk et al.~\cite{Altinoluk:2014eka}:
\begin{itemize}
\item The splitting function has the normalisation of the LO $P_{gg}$, $2N_c[\frac{\xi}{(1-\xi)_+}+\frac{1-\xi}{\xi}+\xi(1-\xi)]+\frac{\beta_0}{2}\delta(1-\xi)$ at $N_f=0$, and the coefficient $\beta_0/2$ of $\delta(1-\xi)$ is the one required by the momentum sum rule $\int_0^1d\xi\,\xi P_{gg}(\xi)=0$. The prefactor $\frac{\alpha_s}{2\pi}$, the factor $\xi^{-2}$ with the phase $e^{-i\v k\cdot(\v y-\v y')/\xi}$, the sign and the dependence on $\mu$ are those of a final-state $\MS$ fragmentation of the LO spectrum of the parent gluon.
\item The end point $\xi\to1$ (integrated gluon soft) is plus-subtracted into the JIMWLK term, as the rapidity divergence is subtracted into BK evolution in~\cite{Chirilli:2012jd}. The end point $\xi\to0$ (observed gluon soft with respect to its parent) is subtracted into the BFKL term; the combination of the two subtractions is exactly the ``double plus'' prescription of~\cite{Kovner:2023vsy}, for the same reason: both poles are part of the LL evolution. In the forward (hybrid) case $\xi$ is bounded from below and no subtraction at $\xi\to0$ arises.
\item The argument of the collinear logarithm is the transverse separation of the observed gluon in the amplitude and in the conjugate amplitude, $(\v y-\v y')/\xi$, whereas in the $\MS$ form of~\cite{Chirilli:2012jd,Altinoluk:2014eka} it is the separation of the parent, $\v y-\v y'$. The difference is a finite term $\int d\xi\,P_{gg}(\xi)\ln\xi^2\otimes$LO, which depends on the $d$-dimensional definition of the convolution. Our form coincides with the cut-off scheme of~\cite{Altinoluk:2014eka}, with $\mu^2_{\MS}=c_0^2\mu^2_{\rm cut}\approx1.26\,\mu^2_{\rm cut}$, $c_0=2e^{-\gamma_E}$. \pending{P4: final form of the logarithm and of the $\xi\to0$ subtraction.}
\item No counterpart arises to the initial-state DGLAP evolution of the projectile parton distribution of the hybrid formalism: the projectile here is described by an unintegrated charge correlator, whose small-$x$ evolution is the BFKL term.
\item The $\beta_0$ terms. The real emission produces no $\delta(1-\xi)$ term, and the virtual $\frac{11}{3}$ terms of the vertex correction (Row I.II) cancel against those of Row I.V. The term $\frac{\beta_0}{2}\delta(1-\xi)$ of $P_{gg}$ is added and subtracted, the subtracted term being the running-coupling term; with it the LO coupling becomes $g^2(\mu)$ at the scale of the collinear subtraction, or $g^2(c_0^2/r^2)$ with $r$ the transverse separation conjugate to $\v k$. In the classification of~\cite{Kovner:2023vsy}, where of the $2\beta_0$ of the NLO JIMWLK kernel one is the ultraviolet running of the coupling and one a DGLAP-type resolution logarithm, our $\beta_0$ terms are of the second, collinear, type. This agrees with Kovchegov and Weigert~\cite{Kovchegov:2007vf} and with Horowitz and Kovchegov~\cite{Horowitz:2010yg}, who find that the coupling of the emission vertex of the produced gluon runs with the collinear resolution scale, $\alpha_s(\Lambda^2_{\rm coll})$, which plays the role of the factorisation scale of the fragmentation function. The other couplings of~\cite{Horowitz:2010yg}, associated with the $t$-channel exchanges, are not renormalised in this calculation. \pending{P3: account of all $1/\epsilon$ poles; wording of the running-coupling term.}
\end{itemize}

\subsection{Finite terms}\label{sec:lit_finite}
% Depends on P1, P2, P5, P6. Revise after sec_finite is filled.

\pending{To be completed after Section~\ref{sec:finite} (P1, P2, P5, P6).} The following properties of the finite terms are established independently of the open points: (i) all finite terms vanish for $U=1$ and for constant $U$; (ii) Group II is real without its complex conjugate; (iii) in Group II, Row I, individual terms have denominators that vanish on whole surfaces in the transverse plane (e.g.\ $(\v x'-\v w')^2(\v x-\v z)^2=(\v x-\v w)^2(\v x'-\v z)^2$), but the sum is finite and continuous across them, so that no principal-value prescription is needed there; this is the mechanism known from the kernels $\ln(X^2Y'^2/X'^2Y^2)/(X^2Y'^2-X'^2Y^2)$ of NLO JIMWLK~\cite{Balitsky:2013fea,Kovner:2013ona}; (iv) the kernels are integrable in all coincidence limits of Groups II and III and of Group I Row V; (v) the collinear limit of Row II.IV agrees with collinear factorisation (Section~\ref{sec:lit_dglap}), and its limits $\xi\to0$ and $\xi\to1$ reproduce the coefficients of the projectile-side and target-side soft-gluon emission, as in the corresponding limits of the NLO jet vertex~\cite{Bartels:2006hg}. The comparison of the full dilute limit with the jet vertex of~\cite{Bartels:2006hg,Fadin:1998py} is structural only. Open at the time of writing are the large-distance behaviour of the soft remainders in Group I Row III and Group III Rows III, V, VI, and the large-distance tail of Row II.IV for a colour-neutral projectile, $\propto\beta_0\frac{4\pi}{\v k^2}\int\frac{d^2Y}{Y^2}[-\frac{11}{6}+\frac16\cos2\varphi]$, for which no compensating term has been identified. The kinematic-constraint remainders $[k^+/(p^+(p^++k^+))]_+$ of Group III are the analogue of the terms discussed in the forward NLO literature~\cite{Stasto:2013cha,Altinoluk:2014eka,Ducloue:2016shw}.

% ======================================================================
\section{Conclusions}\label{sec:conclusions}
% ======================================================================
% SLOT (partially drafted). Inputs: all sections; the paragraphs on the collinear, running-coupling
% and finite terms must be written after sections dglap_rc and finite (P1-P6) are filled in.

We have computed the single inclusive gluon spectrum at mid rapidity in the collision of a dilute projectile with a dense target to next-to-leading order in the Color Glass Condensate, in pure Yang--Mills theory. The calculation uses the evolution operator $\Omega$ of the projectile, written as a normal-ordered series in the soft gluon operators with coefficients that are functionals of the non-commuting valence charge density, and reduces the spectrum to products of these coefficients and of Wilson lines. We have organised the $O(g^4)$ cross section into three groups of terms and fifteen rows.

The rapidity-enhanced terms have the structure required by high-energy factorisation. The logarithm of the longitudinal cut-offs splits at the momentum of the produced gluon. The coefficient of $\ln(\vee/k^+)$ is the BFKL evolution of the projectile charge correlator inserted into the LO spectrum, and the coefficient of $\ln(k^+/\Lambda)$ is the JIMWLK Hamiltonian acting on the target operator of the LO spectrum, for an arbitrary ordered projectile correlator. Terms with four charges cancel, as required by the coherent-state structure of the LO wave function, and the terms with three charges cancel in the symmetrised sense; with non-commuting charges, their reordering produces part of the two-charge structures of both evolutions. The operator ordering of the charges, in particular of the three-charge term of the $O(g^3)$ wave function, is essential for these cancellations.

\pending{To be completed: collinear term and fragmentation, running coupling, finite terms, open points, outlook.}


\acknowledgments
% To be written by the authors.
ACKNOWLEDGMENTS.

\appendix
% ======================================================================
\section{States, Hamiltonian and light-cone matrix elements}\label{app:hamiltonian}
% ======================================================================

In this appendix we collect the definitions of the states, the parts of the eikonal light-cone Hamiltonian that are needed, and the matrix elements used in the LCPT calculations of Appendix~\ref{app:wavefunctions}. They are those of~\cite{Radhakrishnan:2026soft,Lublinsky:2016meo}, restricted to pure Yang--Mills theory.

\subsection{States and normalisation}\label{app:states}

The soft vacuum $\ket0$ is annihilated by all soft annihilation operators, $a^a_i(k^+,\v k)\ket0=0$ for $\Lambda<k^+<\vee$, and its energy is set to zero, $E_0=0$. The one- and two-gluon states are
\begin{align}
\ket{g^a_i(k)}&=\frac{a^{\dagger a}_i(k)}{(2\pi)^{3/2}}\ket0\,,&
\ket{g^a_i(k^+,\v z)}&=\frac{a^{\dagger a}_i(k^+,\v z)}{(2\pi)^{1/2}}\ket0\,,\notag\\
\ket{g^a_i(k)\,g^b_j(p)}&=\frac{a^{\dagger a}_i(k)\,a^{\dagger b}_j(p)}{(2\pi)^3}\ket0\,,&
\ket{g^a_i(k^+,\v z)\,g^b_j(p^+,\v z')}&=\frac{a^{\dagger a}_i(k^+,\v z)\,a^{\dagger b}_j(p^+,\v z')}{2\pi}\ket0\,,
\end{align}
with energies $E_g(k)=\v k^2/2k^+$ and $E_{gg}(k,p)=\v k^2/2k^++\v p^2/2p^+$, and normalisation
\begin{align}
\braket{g^b_j(p)|g^a_i(k)}&=\delta^{ab}\delta_{ij}\,\delta(k^+-p^+)\,\delta^{(2)}(\v k-\v p)\,,\notag\\
\braket{g^b_j(p^+,\v z')|g^a_i(k^+,\v z)}&=\delta^{ab}\delta_{ij}\,\delta(k^+-p^+)\,\delta^{(2)}(\v z-\v z')\,.
\end{align}
The coordinate-space state differs from the Fourier transform of the momentum-space state by a factor $2\pi$, $\ket{g(k^+,\v z)}=2\pi\int\frac{d^2k}{(2\pi)^2}e^{-i\v k\cdot\v z}\ket{g(k)}$. Consequently $a^a_i(k^+,\v k)\ket{g^b_j(q)}=(2\pi)^{3/2}\delta^{ab}\delta_{ij}\delta^{(3)}(k-q)\ket0$.

\subsection{Soft and valence modes}

The gauge field in the light-cone gauge is decomposed into soft and valence components,
\begin{align}
\underline{A}^a_i(x)&=\int_\Lambda^\vee\frac{dk^+}{2\pi}\int\frac{d^2k}{(2\pi)^2}\frac{1}{\sqrt{2k^+}}\Big(a^a_i(k^+,\v k)e^{-ik\cdot x}+a^{\dagger a}_i(k^+,\v k)e^{ik\cdot x}\Big)\,,\notag\\
\overline{A}^a_i(x)&=\int_\vee^\infty\frac{dk^+}{2\pi}\int\frac{d^2k}{(2\pi)^2}\frac{1}{\sqrt{2k^+}}\Big(a^a_i(k^+,\v k)e^{-ik\cdot x}+a^{\dagger a}_i(k^+,\v k)e^{ik\cdot x}\Big)\,.
\end{align}
The valence charge density is the colour charge density of the valence modes~\cite{Radhakrishnan:2026soft},
\begin{align}
\rho^a(-\v p)=-if^{abc}\int_\vee^\infty\frac{dk^+}{2\pi}\int\frac{d^2k}{(2\pi)^2}\,a^{\dagger b}_j(k^+,\v k)\,a^c_j(k^+,\v k+\v p)\,,
\end{align}
and satisfies the algebra~\eqref{eq:rho}. Inserting the decomposition into the light-cone Hamiltonian and dropping all terms suppressed by powers of the ratio of soft to valence longitudinal momenta (in particular all inverse longitudinal derivatives acting on valence fields), one obtains the interaction terms~\cite{Lublinsky:2016meo,Radhakrishnan:2026soft}
\begin{align}
H_g&=-g\int dx^-d^2x\,\big(\partial_i\underline A^a_i\big)\,f^{abc}\,\frac{1}{\partial^+}\big(\overline A^b_j\partial^+\overline A^c_j\big)\,,\notag\\
H_{ggg}&=-gf^{abc}\int dx^-d^2x\,\Big[\big(\partial_i\underline A^a_j\big)\underline A^b_i\underline A^c_j+\big(\partial_i\underline A^a_i\big)\frac{1}{\partial^+}\big(\underline A^b_j\partial^+\underline A^c_j\big)\Big]\,,\notag\\
H_{gg\text{-inst}}&=g^2f^{abc}f^{ade}\int dx^-d^2x\,\frac{1}{\partial^+}\big(\underline A^b_i\partial^+\underline A^c_i\big)\,\frac{1}{\partial^+}\big(\overline A^d_j\partial^+\overline A^e_j\big)\,,\notag\\
H_{gggg}&=\frac{g^2}{4}f^{abc}f^{ade}\int dx^-d^2x\,\Big[\underline A^b_i\underline A^c_j\underline A^d_i\underline A^e_j-2\,\frac{1}{\partial^+}\big(\underline A^b_i\partial^+\underline A^c_i\big)\frac{1}{\partial^+}\big(\underline A^d_j\partial^+\underline A^e_j\big)\Big]\,.
\end{align}
In terms of the soft creation and annihilation operators $H_g$ takes the form~\eqref{eq:Hg}, and
\begin{align}
H_{gg\text{-inst}}=\ &\int_\Lambda^\vee\frac{dp^+}{2\pi}\frac{dq^+}{2\pi}\int\frac{d^2p}{(2\pi)^2}\frac{d^2q}{(2\pi)^2}\frac{ig^2f^{abc}}{\sqrt{p^+q^+}}\notag\\&\times\Bigg(\frac{q^+\rho^a(-\v p-\v q)}{2(p^++q^+)^2}\,a^{\dagger b}_i(p^+,\v p)\,a^{\dagger c}_i(q^+,\v q)-{\rm h.c.}
\notag\\&
-\frac{(p^++q^+)\,\rho^a(-\v p+\v q)}{2(p^+-q^+)^2}\,a^{\dagger b}_i(p^+,\v p)\,a^c_i(q^+,\v q)\Bigg)\,.
\end{align}
The four-gluon interaction $H_{gggg}$ does not enter the coefficients used in this paper.
% CHOICE FOR THE AUTHORS: confirm that H_gggg (normal ordered) gives no O(g^3) one-gluon contribution (no H_gggg diagram appears in the calculation files).

\subsection{Matrix elements}\label{app:matrixelements}

The matrix elements of the interaction terms between the free states, which enter the LCPT expansion~\eqref{eq:lcpt}, are~\cite{Radhakrishnan:2026soft}
\begin{align}
\braket{g^a_i(k)|H_g|0}&=\frac{g\,k^i\rho^a(-\v k)}{4\pi^{3/2}|k^+|^{3/2}}\,,\label{eq:m1}\\
\braket{g^c_l(q)\,g^b_j(p)|H_g|g^a_i(k)}&=\delta^{ac}\delta_{il}\delta^{(3)}(k-q)\frac{g\,p^j\rho^b(-\v p)}{4\pi^{3/2}|p^+|^{3/2}}\notag\\&\quad+\delta^{ab}\delta_{ij}\delta^{(3)}(k-p)\frac{g\,q^l\rho^c(-\v q)}{4\pi^{3/2}|q^+|^{3/2}}\,,\label{eq:m2}\\
\braket{g^c_l(q)\,g^b_j(p)|H_{ggg}|g^a_i(k)}&=\frac{igf^{abc}\,\delta^{(3)}(k-p-q)}{8\pi^{3/2}\sqrt{k^+p^+q^+}}\Bigg[\Big(p^i-q^i+\frac{q^+-p^+}{k^+}k^i\Big)\delta_{jl}\notag\\&\hspace{1cm}+\Big(k^j+q^j-\frac{k^++q^+}{p^+}p^j\Big)\delta_{il}
\notag\\&\hspace{1cm}
+\Big(\frac{k^++p^+}{q^+}q^l-p^l-k^l\Big)\delta_{ij}\Bigg]\,,\label{eq:m3}\\
\braket{g^c_j(q)\,g^b_i(p)|H_{gg\text{-inst}}|0}&
=\frac{ig^2f^{abc}\,(q^+-p^+)\,\delta_{ij}\,\rho^a(-\v p-\v q)}{2(2\pi)^3\sqrt{p^+q^+}\,(p^++q^+)^2}\,,\label{eq:m4}\\
\braket{g^c_j(q)|H_{gg\text{-inst}}|g^b_i(p)}&
=\frac{ig^2f^{abc}\,(p^++q^+)\,\delta_{ij}\,\rho^a(\v p-\v q)}{2(2\pi)^3\sqrt{p^+q^+}\,(p^+-q^+)^2}\,.\label{eq:m5}
\end{align}
Eq.~\eqref{eq:m1} is the emission of a soft gluon by the valence charges, Eq.~\eqref{eq:m2} the emission of a second gluon in the presence of a first one, Eq.~\eqref{eq:m3} the triple-gluon vertex among soft gluons, Eq.~\eqref{eq:m4} the instantaneous emission of two soft gluons by the valence charges and Eq.~\eqref{eq:m5} the instantaneous interaction of a soft gluon with the valence charges, in which the longitudinal momentum $p^+-q^+$ is exchanged with the valence modes. All matrix elements are operators on the valence Hilbert space.

\subsection{Rules of light-cone perturbation theory}\label{app:lcptrules}

For the dressing of the soft vacuum the energy denominators are $E_0-E_n=-E_n$, Eq.~\eqref{eq:lcpt}. For the dressing of a one-gluon state $\ket{g^a_i(k)}$ the energy of the initial state is $E_g(k)$ and the denominators are $E_g(k)-E_n$; for example, the instantaneous transition of the gluon into a gluon of momentum $p$ is
\begin{align}
\ket{g^a_i(k)}\to\int dp^+\!\int d^2p\;\ket{g^b_j(p)}\,\frac{\braket{g^b_j(p)|H_{gg\text{-inst}}|g^a_i(k)}}{E_g(k)-E_g(p)}\,.
\end{align}
These denominators vanish inside the integration region, on $p^+\v k^2=k^+\v p^2$ in momentum space or on $p^+(\v x-\v u)^2=k^+(\v x-\v z)^2$ in coordinate space; they are understood as principal values (Appendix~\ref{app:distributions}). The longitudinal momentum of every soft gluon, real or intermediate, is restricted to $\Lambda<p^+<\vee$, and the longitudinal momentum exchanged with the valence charges in Eqs.~\eqref{eq:m2} and~\eqref{eq:m5} satisfies $|p^+-q^+|>\Lambda$. The kinematic regions allowed for the transition $g(k^+)\to g(p^+)$ through the valence charges are shown in Figure~\ref{fig:kinem}.

\begin{figure}[t]
\centering
\tikzfig{kinem}
\caption{Allowed longitudinal momenta for the transition of a soft gluon of momentum $k^+$ into a soft gluon of momentum $p^+$ by the interaction with the valence charges, for $p^+>k^+$. The upper (purple) region corresponds to $\Lambda<p^+-k^+<\vee$, the green region to $\Lambda<k^+<\vee$ and $\Lambda<p^+<\vee$; the allowed region is their intersection. The case $p^+<k^+$, in which the gluon transfers longitudinal momentum to the valence charges, is the mirror image.}
\label{fig:kinem}
\end{figure}

% ======================================================================
\section{Light-cone wave functions and the coefficients of \texorpdfstring{$\Omega$}{Omega}}\label{app:wavefunctions}
% ======================================================================

The coefficients of $\Omega$ listed in Section~\ref{sec:coefficients} are obtained by computing the action of $\Omega$ on the soft vacuum and on a one-gluon state in two ways: from the expansion~\eqref{eq:OmegaExp}, and from LCPT, Eq.~\eqref{eq:lcpt}, with the matrix elements of Appendix~\ref{app:matrixelements}. The two are equated sector by sector in Fock space. In this appendix we give the Fock components, the diagrams from which they arise, and the steps that lead from the components to the coefficients. In the diagrams the horizontal line denotes the valence charges, curly lines soft gluons, and vertical dashed lines the intermediate states whose energies enter the denominators; a vertical bar on a line denotes an instantaneous interaction. Momenta are written in momentum space; the coordinate forms follow with the Fourier transforms of Appendix~\ref{app:fourier}.

\subsection{Matching conditions}

From Eq.~\eqref{eq:OmegaExp} and the normalisation of the states,
\begin{align}\label{eq:Omegavac}
\Omega\ket0=\mf N\ket0+(2\pi)^{3/2}\int\frac{d^3p}{(2\pi)^3}\Big[\mf A^{(3)b}_{2j}(p)-\mf A^{(1)\dagger b}_{1j}(p)\Big]\ket{g^b_j(p)}+(2\pi)^3\int\frac{d^3p}{(2\pi)^3}\frac{d^3q}{(2\pi)^3}\,\mf B^{cb}_{2kj}(p,q)\ket{g^b_j(p)\,g^c_k(q)}+\dots
\end{align}
and
\begin{align}\label{eq:Omegag}
\Omega\ket{g^a_i(k)}=\ &\frac{1}{(2\pi)^{3/2}}\Big[\mf A^{(1)a}_{1i}(k)+\mf A^{(3)a}_{1i}(k)\Big]\ket0+\mf N\ket{g^a_i(k)}+\int\frac{d^3p}{(2\pi)^3}\,\mf B^{ba}_{3ji}(k,p)\ket{g^b_j(p)}
\notag\\&
+(2\pi)^{3/2}\int\frac{d^3p}{(2\pi)^3}\Big[\mf A^{(3)b}_{2j}(p)-\mf A^{(1)\dagger b}_{1j}(p)\Big]\ket{g^a_i(k)\,g^b_j(p)}
+(2\pi)^{3/2}\int\frac{d^3p}{(2\pi)^3}\,\mf C^{bca}_{1jki}(k,p)\ket{g^b_j(p)\,g^c_k(k-p)}+\dots,
\end{align}
where the dots denote states with more gluons. The vacuum sector fixes $\mf N$, $\mf A^{(1)}$, $\mf A^{(3)}_2$ and $\mf B_2$; the one-gluon sector fixes $\mf B_3$ and $\mf C_1$, and reproduces $\mf A^{(1)}_1$ and the product $\mf A\,$(one-gluon state) as a consistency check.

\subsection{The soft vacuum at order \texorpdfstring{$g$}{g} and the normalisation}

At order $g$ the valence charges emit one soft gluon (Figure~\ref{fig:vac_g}):
\begin{align}\label{eq:PsiLO}
\ket{\Psi^{\rm LO}}=-\int_\Lambda^\vee dk^+\!\int d^2k\,\ket{g^a_i(k)}\frac{\braket{g^a_i(k)|H_g|0}}{E_g(k)}
=-\int_\Lambda^\vee dk^+\!\int d^2k\,\frac{g\,k^i\rho^a(-\v k)}{2\pi^{3/2}\sqrt{k^+}\,\v k^2}\ket{g^a_i(k)}
=\frac{ig}{2\pi^{3/2}}\int_\Lambda^\vee\frac{dk^+}{\sqrt{k^+}}\int_{x,z}\Kv^i(\v x-\v z)\rho^a(\v x)\ket{g^a_i(k^+,\v z)}\,.
\end{align}
Comparison with Eq.~\eqref{eq:Omegavac} gives $-\mf A^{(1)a}_{2i}(k^+,\v k)=\mf A^{(1)\dagger a}_{1i}(k^+,\v k)=\sqrt2\,g\,k^i\rho^a(-\v k)/(\sqrt{k^+}\v k^2)$, i.e.\ Eq.~\eqref{eq:A1}. The norm of the state to order $g^2$ fixes
\begin{align}
\mf N=1-\frac12\braket{\Psi^{\rm LO}|\Psi^{\rm LO}}=1-\frac{g^2}{8\pi^3}\int_\Lambda^\vee\frac{dp^+}{p^+}\int d^2p\,\frac{\rho^a(-\v p)\rho^a(\v p)}{\v p^2}\,,
\end{align}
which is Eq.~\eqref{eq:N2}.

\begin{figure}[t]
\centering
\tikzfig{vacuumtog}
\caption{Emission of a soft gluon by the valence charges at order $g$.}
\label{fig:vac_g}
\end{figure}

\subsection{Two soft gluons at order \texorpdfstring{$g^2$}{g2}}

Three diagrams contribute (Figure~\ref{fig:vac_gg}): the emission of two gluons by the valence charges, with the matrix elements~\eqref{eq:m1} and~\eqref{eq:m2},
\begin{align}
\ket{\Psi^{1}_{gg}}=\frac12\int dq^+dk^+dp^+\!\int d^2q\,d^2k\,d^2p\,\ket{g^b_j(k-p)\,g^c_l(p)}\frac{\braket{g^b_j(k-p)\,g^c_l(p)|H_g|g^a_i(q)}\braket{g^a_i(q)|H_g|0}}{E_{gg}(p,k-p)\,E_g(q)}\,,
\end{align}
the emission of one gluon which splits into two through the triple-gluon vertex~\eqref{eq:m3},
\begin{align}
\ket{\Psi^{2}_{gg}}=\frac12\int dk^+dp^+dq^+\!\int d^2k\,d^2p\,d^2q\,\ket{g^b_j(q)\,g^c_l(p)}\frac{\braket{g^b_j(q)\,g^c_l(p)|H_{ggg}|g^a_i(k)}\braket{g^a_i(k)|H_g|0}}{E_{gg}(p,q)\,E_g(k)}\,,
\end{align}
and the instantaneous emission of the pair~\eqref{eq:m4},
\begin{align}
\ket{\Psi^{3}_{gg}}=-\frac12\int dk^+dp^+\!\int d^2k\,d^2p\,\ket{g^b_i(k-p)\,g^c_j(p)}\frac{\braket{g^b_i(k-p)\,g^c_j(p)|H_{gg\text{-inst}}|0}}{E_{gg}(k-p,p)}\,.
\end{align}
In the first diagram the product of the two charges is decomposed as
\begin{align}\label{eq:decomp}
\rho^b(-\v k+\v p)\,\rho^a(-\v p)=\frac{i}{2}f^{bad}\rho^d(-\v k)+\frac12\big\{\rho^b(-\v k+\v p),\rho^a(-\v p)\big\}\,.
\end{align}
The commutator term is linear in $\rho$ and is combined with the other two diagrams. The anticommutator term describes two independent emissions. Comparing with Eq.~\eqref{eq:Omegavac} gives the coefficient
\begin{align}\label{eq:B2mom}
    \mf{B}^{ba}_{2ji}(k^{+},k;p^{+},p) & =  \frac{ig^{2}f^{abc}\,p^{j}k^{i}\rho^{c}(-p-k)\, \sqrt{p^{+}k^{+}}}{2\,\left(p^{+}k^{2}+k^{+}p^{2}\right)}\,\left[\frac{1}{k^{+}\,p^{2}} - \frac{1}{k^{2}\,p^{+}}\right]\,\Theta\left(\vee-k^{+}-p^{+}\right)\notag \\ & + \frac{g^{2}p^{j}k^{i}\,\{\rho^{b}(-p),\rho^{a}(-k)\}}{2\,p^{2}k^{2}\, \sqrt{p^{+}k^{+}}}\,\Theta\left(\vee-k^{+}\right)\Theta\left(\vee-p^{+}\right)\notag \\ & -\,\frac{ig^{2}f^{abc}\rho^{c}(-p-k)\,(p+k)^{l}\,\left(p^{m}k^{+}-p^{+}k^{m}\right)}{(p+k)^{2}\,\sqrt{p^{+}k^{+}}\,(p^{+}+k^{+})\,\left(\frac{p^{2}}{p^{+}}+\frac{k^{2}}{k^{+}}\right)}\,\left[\frac{\delta_{lm}\delta_{ji}}{(p^{+}+k^{+})} - \frac{\delta_{im}\delta_{jl}}{k^{+}} - \frac{\delta_{jm}\delta_{il}}{p^{+}}\right]\notag \\ & \qquad\times\Theta\left(\vee-k^{+}-p^{+}\right)\notag \\ & -\,\frac{ig^{2}f^{abc}\rho^{c}(-p-k)\,(p^{+}-k^{+})\delta_{ji}}{2\,\sqrt{p^{+}k^{+}}\,(p^{+}+k^{+})^{2}\,\left(\frac{k^{2}}{k^{+}}+\frac{p^{2}}{p^{+}}\right)}\,\Theta\left(\vee-k^{+}-p^{+}\right)
\end{align} [$\Theta$ replaced; the support follows from the kinematics ($k^+,p^+>\Lambda$ throughout). Third term (triple vertex): the parent gluon carries $k^++p^+$ and is itself a soft mode, hence $\Theta(\vee-k^+-p^+)$; the instantaneous fourth term has the same support in Eq.~\eqref{ggthree} (a convention; no particle carries $k^++p^+$ there). Second term ($\{\rho,\rho\}$): the two gluons are emitted independently by the valence charges, so $\Lambda<k^+<\vee$ and $\Lambda<p^+<\vee$ with no condition on the sum, as for the same intermediate state in Eq.~\eqref{3rho}; with $k^++p^+<\vee$ the four-$\rho$ terms of the cross section would not cancel for commuting $\rho$ (coherent state), see the note after Eq.~\eqref{eq:CS1}. The first term is the commutator part of the same diagram, so by the same kinematics its support is also the full square; it is left unchanged here for the author to decide. $\Theta(k^+-\Lambda-p^+)$ is not symmetric under exchange of the two gluons and removes $p^+>k^+$.]\par\noindent 
where the first, third and fourth terms come from the commutator part of the first diagram, from the triple-gluon vertex and from the instantaneous emission, respectively, and the second term from the anticommutator part. The supports are discussed in Section~\ref{sec:coefficients}. With the Fourier transforms~\eqref{eq:FT_i2} and~\eqref{eq:FT_i7} this becomes Eq.~\eqref{eq:B2}.

\begin{figure}[t]
\centering
\tikzfig{vacuumto2g2rho}\qquad\tikzfig{vacuumto2g}\qquad\tikzfig{vacuumto2ginst}
\caption{Two soft gluons at order $g^2$: two emissions by the valence charges, emission followed by splitting, and instantaneous emission.}
\label{fig:vac_gg}
\end{figure}

\subsection{One soft gluon at order \texorpdfstring{$g^3$}{g3}}

\paragraph{Diagrams.} The one-gluon component at order $g^3$ receives contributions linear, quadratic and cubic in the charges.
\begin{itemize}
\item The self-energy of the emitted gluon through two triple-gluon vertices (Figure~\ref{fig:vac_g3_onerho}),
\begin{align}
\ket{\Psi^{1}_{g\rho}}=-\frac12\int dk^+dp^+dq^+dr^+\!\int d^2k\,d^2p\,d^2q\,d^2r\,\ket{g^d_l(r)}\frac{\braket{g^d_l(r)|H_{ggg}|g^c_k(q)g^b_j(p)}\braket{g^c_k(q)g^b_j(p)|H_{ggg}|g^a_i(k)}\braket{g^a_i(k)|H_g|0}}{E_g(r)\,E_{gg}(p,q)\,E_g(k)}\,.
\end{align}
It is linear in $\rho$ and is the only diagram with an ultraviolet divergence of the transverse loop integral.
\item The diagrams with two valence vertices (Figure~\ref{fig:vac_g3_tworho}): the emitted gluon interacts instantaneously with the valence charges, gaining or losing longitudinal momentum; two gluons emitted by the valence charges recombine through the triple-gluon vertex; two gluons are emitted, one instantaneously, and one of them is absorbed by the valence charges; a gluon splits into two, one of which is absorbed by the valence charges. With the decomposition~\eqref{eq:decomp} each of them gives a term linear in $\rho$ (commutator) and a term $\propto f^{abc}\{\rho^b,\rho^c\}$.
\item The diagrams with three valence vertices (Figure~\ref{fig:vac_g3_threerho}): a gluon is emitted and reabsorbed by the valence charges before or after the emission of the observed gluon, in the three time orderings, and the corresponding normalisation term of Eq.~\eqref{eq:lcpt}.
\end{itemize}

\begin{figure}[t]
\centering
\tikzfig{loop}
\caption{Self-energy of the emitted gluon at order $g^3$.}
\label{fig:vac_g3_onerho}
\end{figure}

\begin{figure}[t]
\centering
\tikzfig{Misha_2rho_single_g}\qquad\tikzfig{Misha_Out_2rho_single_g}\qquad\tikzfig{Misha_2rho_3_single_g}\\[4mm]
\tikzfig{Misha_2rho_insta_single_g}\qquad\tikzfig{Misha_2rho_loop_single_g}
\caption{Diagrams with two valence vertices contributing to the one-gluon component at order $g^3$.}
\label{fig:vac_g3_tworho}
\end{figure}

\begin{figure}[t]
\centering
\tikzfig{Misha_2rho_loopVal1_single_g}\qquad\tikzfig{Misha_2rho_loopVal2_single_g}\qquad\tikzfig{Misha_2rho_loopVal3_single_g}
\caption{Diagrams with three valence vertices contributing to the one-gluon component at order $g^3$.}
\label{fig:vac_g3_threerho}
\end{figure}

\paragraph{Self-energy.} After the transverse loop integral in $d=2-\epsilon$ dimensions, with $p^+=\xi k^+$,
\begin{align}
\ket{\Psi^{1}_{g\rho}}=-\int_\Lambda^\vee\frac{dk^+}{\sqrt{k^+}}\int d^2k\int_{\Lambda/k^+}^{1-\Lambda/k^+}d\xi\,\frac{g^3N_c\,\rho^a(-\v k)\,k^i}{16\pi^{7/2}\,\v k^2}\Big[-2+\xi(1-\xi)+\frac{1}{\xi(1-\xi)}\Big]\Big[-\frac2\epsilon+\ln\frac{\v k^2\xi(1-\xi)}{\mu^2_{\MS}}\Big]\ket{g^a_i(k^+,\v k)}\,.
\end{align}
With $\delta=\Lambda/k^+\to0$,
\begin{align}\label{eq:xiintegrals}
\int_\delta^{1-\delta}d\xi\Big[-2+\xi(1-\xi)+\frac{1}{\xi(1-\xi)}\Big]=-\frac{11}{6}-2\ln\delta\,,\qquad
\int_0^1d\xi\,\big[-2+\xi(1-\xi)\big]\ln\big(\xi(1-\xi)\big)=\frac{67}{18}\,,\qquad
\int_\delta^{1-\delta}d\xi\,\frac{\ln(\xi(1-\xi))}{\xi(1-\xi)}=-\ln^2\delta-\frac{\pi^2}{3}\,,
\end{align}
so that
\begin{align}
\ket{\Psi^{1}_{g\rho}}=\int_\Lambda^\vee\frac{dk^+}{\sqrt{k^+}}\int d^2k\,\frac{g^3N_c\,\rho^a(-\v k)\,k^i}{32\pi^{7/2}\,\v k^2}\Big[\Big(\frac{11}{3}+4\ln\frac{\Lambda}{k^+}\Big)\Big(-\frac2\epsilon+\ln\frac{\v k^2}{\mu^2_{\MS}}\Big)+2\ln^2\frac{\Lambda}{k^+}-\frac{67}{9}+\frac{2\pi^2}{3}\Big]\ket{g^a_i(k^+,\v k)}\,.
\end{align}
This is the origin of the coefficient $\frac{11}{3}=\beta_0/N_c$ and of the constant $\frac{67}{9}$.

\paragraph{Sum of the one-$\rho$ terms.} Adding the commutator parts of the diagrams with two and three valence vertices, the one-gluon component linear in $\rho$ is
\begin{align}\label{eq:Psigrho}
\ket{\Psi_{g\rho}}=\ &\int_\Lambda^\vee\frac{dk^+}{\sqrt{k^+}}\int d^2k\,\frac{g^3N_c\,\rho^a(-\v k)\,k^i}{32\pi^{7/2}\,\v k^2}\Big[\Big(-\frac2\epsilon+\ln\frac{\v k^2}{\mu^2_{\MS}}\Big)\Big(\frac{11}{3}+\ln\frac{\Lambda}{k^+}-2\ln\frac{\vee}{\Lambda}\Big)+\frac{\pi^2}{3}-\frac{67}{9}+\ln^2\frac{\Lambda}{k^+}-\ln^2\frac{\vee}{k^+}\Big]\ket{g^a_i(k^+,\v k)}
\notag\\&
+\int_\Lambda^\vee\frac{dk^+}{\sqrt{k^+}}\int d^2k\int_{\Lambda/k^+}^{1-\Lambda/k^+}d\xi\,\frac{g^3N_c\,\rho^a(-\v k)\,k^i}{32\pi^{7/2}\sqrt\xi(1-\xi)\,\v k^2}\Big[(\xi-2)\ln\xi+(\xi+1)\ln(1-\xi)-(\xi-2)\Big(-\frac2\epsilon+\ln\frac{\v k^2}{\mu^2_{\MS}}\Big)\Big]\ket{g^a_i(\xi k^+,\v k)}\,.
\end{align}
The second line, in which the observed gluon carries the momentum fraction $\xi$ of the parent, comes from the diagrams in which a gluon of momentum $k^+$ transfers part of its longitudinal momentum to the valence charges. To read off the coefficient of a gluon of momentum $k^+$ we relabel $\xi k^+\to k^+$ in this term; the parent momentum is then $k^+/\xi$, and the conditions $\Lambda<k^+/\xi<\vee$, $\xi<1-\Lambda\xi/k^+$ give $k^+/\vee<\xi<k^+/(k^++\Lambda)$. The $\xi$ integrals,
\begin{align}
\int_{k^+/\vee}^{k^+/(k^++\Lambda)}d\xi\,\frac{2-\xi}{\xi(1-\xi)}&=2\ln\frac{\vee}{k^+}+\ln\frac{k^+}{\Lambda}+\ln\Big(1-\frac{k^+}{\vee}\Big)+O\Big(\frac{\Lambda}{k^+}\Big)\,,\notag\\
\int_{k^+/\vee}^{k^+/(k^++\Lambda)}d\xi\,\frac{(\xi-2)\ln\xi+(\xi+1)\ln(1-\xi)}{\xi(1-\xi)}&=-\ln^2\frac{\Lambda}{k^+}+\ln^2\frac{\vee}{k^+}+\ln^2\Big(1-\frac{k^+}{\vee}\Big)+\ln\Big(1-\frac{k^+}{\vee}\Big)\ln\frac{\vee}{k^+}+O\Big(\frac{\Lambda}{k^+}\ln\frac{\Lambda}{k^+}\Big)\,,
\end{align}
add $\ln(\vee/k^+-1)+\ln(\vee/\Lambda)$ to the coefficient of the pole and convert the constant of the first line into $\frac{\pi^2}{3}-\frac{67}{9}+\ln^2(1-k^+/\vee)+\ln(1-k^+/\vee)\ln(\vee/k^+)$. The double logarithms of the cut-offs cancel, and the result is Eq.~\eqref{eq:A3onerho}.

\paragraph{Two-$\rho$ terms.} The anticommutator parts of the diagrams with two and three valence vertices give
\begin{align}\label{eq:Psigrhorho}
\ket{\Psi_{g\rho\rho}}=\int_\Lambda^\vee\frac{dk^+}{\sqrt{k^+}}\int d^2k\,d^2p\,\frac{ig^3f^{abc}\{\rho^b(-\v p),\rho^c(-\v k+\v p)\}}{32\pi^{9/2}\,\v p^2\,\v k^2}\Bigg[&\ln\frac{\Lambda}{k^+}\,\frac{\v k^2p^i-(\v k\cdot\v p)k^i}{(\v k-\v p)^2}-2k^i\ln\frac{\vee}{\Lambda}-2k^i\ln\frac{\Lambda\v k^2+k^+\v p^2}{\vee\v k^2+k^+\v p^2}
\notag\\&
+\frac{(2\v p^2-2\v k\cdot\v p+\v k^2)k^i+\v k^2p^i}{(\v k-\v p)^2}\ln\Big[1+\Big(\frac{\vee}{k^+}-1\Big)\frac{\v k^2}{(\v k-\v p)^2}\Big]
\notag\\&
-\frac{\v p^2k^i+\v k^2p^i}{(\v k-\v p)^2}\ln\frac{\vee-k^+}{\Lambda}-k^i\ln\frac{\vee}{k^+}\Bigg]\ket{g^a_i(k)}\,,
\end{align}
after the $\xi$ integration. The terms proportional to $k^+/\Lambda$ of the individual instantaneous diagrams cancel in the sum. Setting $\Lambda\to0$ where possible and transforming to coordinate space gives Eq.~\eqref{eq:A3tworho}.

\paragraph{Three-$\rho$ terms.} The diagrams with three valence vertices give, in addition to commutator terms,
\begin{align}
\ket{\Psi^{6c}_{g\rho\rho\rho}}&=\int_\Lambda^\vee dk^+\!\int d^2k\,d^2p\,\frac{g^3\,k^i\,\rho^b(\v p)\rho^b(-\v p)\rho^a(-\v k)\sqrt{k^+}}{8\pi^{9/2}\,\v k^4}\Big[\frac1\vee-\frac1\Lambda\Big]\ket{g^a_i(k^+,\v k)}\,,\notag\\
\ket{\Psi^{7}_{g\rho\rho\rho}}&=\int_\Lambda^\vee\frac{dk^+}{\sqrt{k^+}}\int d^2k\,d^2p\,\frac{g^3\,k^i\,\rho^a(-\v k)\rho^b(\v p)\rho^b(-\v p)}{16\pi^{9/2}\,\v k^2}\Big[\frac{1}{\v p^2}\ln\frac{\vee}{\Lambda}+\frac{2k^+}{\v k^2}\Big(\frac1\Lambda-\frac1\vee\Big)\Big]\ket{g^a_i(k^+,\v k)}\,.
\end{align}
The second line is the normalisation term of Eq.~\eqref{eq:lcpt}, with the single-emission matrix element to the left of the norm factor. The terms proportional to $1/\Lambda$ and $1/\vee$ cancel exactly between the two lines, although the charges stand in different orders: $\int d^2p\,\rho^b(\v p)\rho^b(-\v p)$ without the weight $1/\v p^2$ is local in coordinate space and commutes with $\rho^a$. What remains is
\begin{align}\label{eq:Psigrhorhorho}
\ket{\Psi_{g\rho\rho\rho}}=\ln\frac{\vee}{\Lambda}\int_\Lambda^\vee\frac{dk^+}{\sqrt{k^+}}\int d^2k\,d^2p\,\frac{g^3\,\rho^a(-\v k)\,\rho^b(\v p)\,\rho^b(-\v p)\,k^i}{16\pi^{9/2}\,\v k^2\,\v p^2}\ket{g^a_i(k^+,\v k)}\,,
\end{align}
which gives $[\mf A^{(3)}_2]_{\rho\rho\rho}=\mf A^{(1)}_2\mf N_2$, Eq.~\eqref{eq:A3threerho}.

\subsection{One soft gluon in the initial state}

The coefficients $\mf B_3$ and $\mf C_1$ are fixed by the dressing of a one-gluon state $\ket{g^a_i(k)}$ at order $g^2$ and $g$, respectively. The diagrams for the transition into a one-gluon state are shown in Figure~\ref{fig:g_g}: the instantaneous interaction with the valence charges, for $p^+>k^+$ and $p^+<k^+$; the emission of a second gluon by the valence charges and the absorption of the first by the triple-gluon vertex, or vice versa; and the diagrams with two valence vertices, in which the first gluon is absorbed and a second one emitted, or the second one emitted before the first is absorbed. The splitting of the gluon into two soft gluons is shown in Figure~\ref{fig:g_gg}. The energy denominators are $E_g(k)-E_n$ (Appendix~\ref{app:lcptrules}). Comparison with Eq.~\eqref{eq:Omegag} gives
\begin{align}\label{eq:B3mom}
    \mf{B}^{ba}_{3ji}(k^{+},k; p^{+},p) & =  \frac{ig^2 f^{abc}\,\sqrt{p^{+}k^{+}} (k^++p^+) } {(p^+-k^+)^2} \frac{\rho^c(-\v{p}+\v{k})}{p^{+}\v{k}^2 - k^{+}\v{p}^2} \delta_{ij}\left[\Theta\left(p^{+} - k^{+} - \Lambda\right)+\Theta\left(k^{+}-\Lambda-p^{+}\right)\right]\notag \\ & -\frac{g^{2}\,k^{i}p^{j}}{k^{2}p^{2}\,\sqrt{k^{+}p^{+}}}\left[\{\rho^{a}(k),\rho^{b}(-p)\}-if^{abc}\rho^{c}(k-p)\frac{p^{2}k^{+}+k^{2}p^{+}}{p^{2}k^{+}-k^{2}p^{+}}\right]\notag \\ & +\frac{2i\,g^{2}f^{abc}\,\rho^{c}(k-p)(k-p)^{l}\,\sqrt{k^{+}p^{+}}}{(k^{+}-p^{+})^{2}}\times \frac{\left(k^{m}p^{+}-k^{+}p^{m}\right)}{\left(k^{2}p^{+}-p^{2}k^{+}\right)}\,\left[\frac{\delta_{jm}\delta_{il}}{p^{+}}-\frac{\delta_{im}\delta_{jl}}{k^{+}}-\frac{\delta_{lm}\delta_{ij}}{(p^{+}-k^{+})}\right]\times \notag \\ & \left[\frac{(p^{+}-k^{+})}{(k-p)^{2}}\,\Theta(p^{+}-k^{+}-\Lambda)+\frac{1}{\left(\frac{k^{2}}{k^{+}}-\frac{(k-p)^{2}}{(k^{+}-p^{+})}-\frac{p^{2}}{p^{+}}\right)}\,\Theta(k^{+}-\Lambda-p^{+})\right]
\end{align} [Sign of the non-instantaneous $p^+>k^+$ term fixed (Eqs.~\eqref{onegrho12} and \eqref{Eq:onegrhoX1} carry the wrong overall sign): it follows from \eqref{m2}, \eqref{m3}, from the cancellation of the $1/(1-\xi)^2$ singularity, and from the $a^\dagger a$ unitarity constraint.]\par\noindent 
and
\begin{align}\label{eq:C1mom}
    \mf{C}^{bca}_{1jki}(k^{+},k; p^{+},p) & = -\,\frac{igf^{abc}\, \left(p^{m}k^{+}-p^{+}k^{m}\right)\,\sqrt{k^{+}p^{+}(k^{+}-p^{+})}}{\sqrt{2}\,\left(pk^{+}-p^{+}k\right)^{2}}\,\left[\frac{\delta_{im}\delta_{jk}}{k^{+}}-\frac{\delta_{ik}\delta_{mj}}{p^{+}}-\frac{\delta_{km}\delta_{ij}}{(k^{+}-p^{+})}\right]\,\Theta\left(k^{+}-\Lambda-p^{+}\right) 
\end{align}
The double poles $1/(k^+-p^+)^2$ of the instantaneous term and of the term with the triple-gluon vertex cancel each other at $|p^+-k^+|\to\Lambda$, so that no term proportional to $k^+/\Lambda$ arises. The coordinate forms, obtained with the Fourier transforms~\eqref{eq:FT_i3}--\eqref{eq:FT_i5}, are Eqs.~\eqref{eq:B3} and~\eqref{eq:C1}.

\begin{figure}[t]
\centering
\tikzfig{gtog_inst}\qquad\tikzfig{gandgtog}\qquad\tikzfig{gto3g_torho}\\[4mm]
\tikzfig{gtovactog_tworho}\qquad\tikzfig{gg_tworho}
\caption{Transition of a soft gluon into a soft gluon through the interaction with the valence charges at order $g^2$.}
\label{fig:g_g}
\end{figure}

\begin{figure}[t]
\centering
\tikzfig{g_togg}\qquad\tikzfig{gtog_plusg}
\caption{Splitting of a soft gluon into two soft gluons ($\rho$-independent, order $g$), and emission of a second soft gluon by the valence charges in the presence of the first one.}
\label{fig:g_gg}
\end{figure}

% ======================================================================
\section{Unitarity relations among the coefficients of \texorpdfstring{$\Omega$}{Omega}}\label{app:unitarity}
% ======================================================================

The coefficients of $\Omega$ that multiply annihilation operators are fixed by the unitarity condition $\Omega^\dagger\Omega=1$, imposed order by order in $g$. With the expansion~\eqref{eq:OmegaExp} and its conjugate, we normal order $\Omega^\dagger\Omega$ with the commutator~\eqref{eq:aadagger}, keeping only the terms needed for the single inclusive cross section at $O(g^4)$ and using $\mf N=\mf N^\dagger$. Since the coefficients are operators on the valence space, their order in each product is kept. The coefficients of the operator structures $1$, $a^\dagger$, $a^\dagger a^\dagger$ and $a^\dagger a^\dagger a$ must vanish (except for the identity), which gives
\begin{align}
\mf N&=1-\frac12\int\frac{d^3p}{(2\pi)^3}\,\mf A^{(1)b}_{1j}(p)\,\mf A^{(1)\dagger b}_{1j}(p)\,,\label{eq:Nunit}\\
\mf A^{(1)\dagger b}_{1j}(p)&=-\mf A^{(1)b}_{2j}(p)\,,\label{eq:A11unit}\\
\mf A^{(3)\dagger b}_{1j}(p)&=-\mf A^{(3)b}_{2j}(p)-\big[\mf N_2,\mf A^{(1)b}_{2j}(p)\big]
+\int\frac{d^3q}{(2\pi)^3}\Big[2\,\mf A^{(1)c}_{1k}(q)\,\mf B^{cb}_{2kj}(p,q)-\mf B^{\dagger cb}_{3kj}(p,q)\,\mf A^{(1)c}_{2k}(q)
\notag\\&\hspace{5.1cm}
-2\,\mf C^{\dagger dcb}_{1lkj}(p,q)\,\mf B^{cd}_{2kl}(q,p-q)\Big]\,,\label{eq:A31unit}\\
\mf B^{\dagger cb}_{1kj}(p,q)&=-\mf B^{cb}_{2kj}(p,q)+\frac12\Big[\mf A^{(1)b}_{2j}(p)\,\mf A^{(1)c}_{2k}(q)+\mf A^{(1)c}_{2k}(q)\,\mf A^{(1)b}_{2j}(p)\Big]+\mf C^{cbd}_{1kjl}(q+p,q)\,\mf A^{(1)d}_{2l}(q+p)\,,\label{eq:B1unit}\\
\mf C^{\dagger dcb}_{2lkj}(p,q)&=-\mf C^{dcb}_{1lkj}(p,q)\,.\label{eq:C2unit}
\end{align}
In addition, the coefficients obey the symmetry relations $\mf B^{cb}_{1kj}(p,q)=\mf B^{bc}_{1jk}(q,p)$, $\mf B^{cb}_{2kj}(p,q)=\mf B^{bc}_{2jk}(q,p)$, $\mf C^{dcb}_{1lkj}(p,q)=\mf C^{cdb}_{1klj}(p,p-q)$, and similarly for $\mf C_2$, which follow from the symmetry of the products of creation and annihilation operators. The coefficient of $a^\dagger a$ gives a relation between $\mf B_3$, $\mf B_3^\dagger$, $\mf A^{(1)}$ and $\mf C_1$, which is not needed for the cross section but serves as a check of $\mf B_3$; in particular it fixes the sign of the non-instantaneous term of $\mf B_3$ for $p^+>k^+$.

In transverse coordinate space the relations read, with the Fourier conventions~\eqref{eq:FTconventions},
\begin{align}
\mf A^{(1)\dagger b}_{1j}(p^+,\v x)&=-\mf A^{(1)b}_{2j}(p^+,-\v x)\,,\notag\\
\mf A^{(3)\dagger b}_{1j}(p^+,\v x)&=-\mf A^{(3)b}_{2j}(p^+,-\v x)-\big[\mf N_2,\mf A^{(1)b}_{2j}(p^+,-\v x)\big]
+\int\frac{dq^+}{2\pi}\int_y\Big[2\,\mf A^{(1)c}_{1k}(q^+,\v y)\,\mf B^{cb}_{2kj}(p^+,-\v x;q^+,-\v y)
\notag\\&\qquad
-\mf B^{\dagger cb}_{3kj}(p^+,\v x;q^+,-\v y)\,\mf A^{(1)c}_{2k}(q^+,-\v y)
-2\int_z\mf C^{\dagger dcb}_{1lkj}(p^+,\v x-\v z;q^+,\v z-\v y)\,\mf B^{dc}_{2lk}(p^+-q^+,-\v z;q^+,-\v y)\Big]\,,\notag\\
\mf B^{\dagger cb}_{1kj}(p^+,\v x;q^+,\v y)&=-\mf B^{cb}_{2kj}(p^+,-\v x;q^+,-\v y)+\frac12\Big[\mf A^{(1)b}_{2j}(p^+,-\v x)\mf A^{(1)c}_{2k}(q^+,-\v y)+\mf A^{(1)c}_{2k}(q^+,-\v y)\mf A^{(1)b}_{2j}(p^+,-\v x)\Big]
\notag\\&\quad
+\int_z\mf C^{cbd}_{1kjl}(q^++p^+,\v z-\v x;q^+,\v x-\v y)\,\mf A^{(1)d}_{2l}(q^++p^+,-\v z)\,,\notag\\
\mf C^{\dagger dcb}_{2lkj}(p^+,\v z;q^+,\v x)&=-\mf C^{dcb}_{1lkj}(p^+,-\v z;q^+,-\v x)\,.
\end{align}

\paragraph{The commutator in $\mf A^{(3)}_1$.} The term $-[\mf N_2,\mf A^{(1)}_2]$ in Eq.~\eqref{eq:A31unit} is absent for commuting charges. It does not vanish here: $\mf N_2\propto\int\rho^a(\v x)\rho^a(\v y)K(\v x,\v y,\v z)$, and $[\rho^a(\v x)\rho^a(\v y),\rho^b(\v w)]\neq0$ for $\v x\neq\v y$, so that $[\mf N_2,\mf A^{(1)}_2]\propto g^3f^{abc}\{\rho^a,\rho^c\}\neq0$. With the ordering $[\mf A^{(3)}_2]_{\rho\rho\rho}=\mf A^{(1)}_2\mf N_2$ derived in Appendix~\ref{app:wavefunctions}, the first two terms of Eq.~\eqref{eq:A31unit} combine to
\begin{align}
-\big[\mf A^{(3)}_2\big]_{\rho\rho\rho}-\big[\mf N_2,\mf A^{(1)}_2\big]=-\mf N_2\,\mf A^{(1)}_2\,.
\end{align}
In the one-gluon amplitude~\eqref{eq:A3amp} this produces the combination~\eqref{eq:Nterms} of normalisation terms, which vanishes for $U=1$, as required by $\bar\Omega^\dagger\Omega=1$ at $U=1$. With the opposite ordering of the three-$\rho$ term, $\mf N_2\mf A^{(1)}_2$, the coefficient $\mf A^{(3)}_2$ would differ by $[\mf N_2,\mf A^{(1)}_2]$, a two-$\rho$ term proportional to $\ln(\vee/\Lambda)$ equal and opposite to the term $-2k^i\ln(\vee/\Lambda)$ of Eq.~\eqref{eq:A3tworho}, and the cross section would acquire a spurious two-$\rho$ term proportional to $\ln(\vee/\Lambda)$ (which vanishes at $U=1$ but not otherwise).

% ======================================================================
\section{Fourier transforms}\label{app:fourier}
% ======================================================================

With the conventions~\eqref{eq:FTconventions}, $f(\v k)=\int_xe^{i\v k\cdot\v x}f(\v x)$ and $f(\v x)=\int\frac{d^2k}{(2\pi)^2}e^{-i\v k\cdot\v x}f(\v k)$, the transverse integrals needed to transform the wave functions and coefficients from momentum to coordinate space are~\cite{Radhakrishnan:2026soft}
\begin{align}
&\int\frac{d^2k}{2\pi}\,\frac{k^i}{\v k^2}\,e^{i\v k\cdot\v X}=\frac{iX^i}{\v X^2}\,,\label{eq:FT_K}\\
&\int\frac{d^2k}{2\pi}\frac{d^2p}{2\pi}\,\frac{k^ip^j}{\v k^2\,(a\v k^2+b\v p^2)}\,e^{i\v k\cdot\v X+i\v p\cdot\v Z}=-\frac{X^iZ^j}{\v Z^2\,(b\v X^2+a\v Z^2)}\,,\label{eq:FT_i2}\\
&\int\frac{d^2k}{2\pi}\frac{d^2p}{2\pi}\,\frac{1}{\xi\v k^2-\v p^2}\,e^{i\v k\cdot\v X+i\v p\cdot\v Z}=\frac{1}{\xi\v Z^2-\v X^2}\,,\label{eq:FT_i3}\\
&\int\frac{d^2k}{2\pi}\frac{d^2p}{2\pi}\,\frac{k^m-p^m}{(\v k-\v p)^2}\,\frac{p^l-\xi k^l}{\v p^2-\xi\v k^2}\,e^{i\v k\cdot\v X+i\v p\cdot\v Y}=-\frac{X^m+\xi Y^m}{\v X^2-\xi\v Y^2}\,\frac{X^l+Y^l}{(\v X+\v Y)^2}\,,\label{eq:FT_i4}\\
&\int\frac{d^2k}{2\pi}\frac{d^2p}{2\pi}\,\frac{k^m-p^m}{\v p^2-\xi\v k^2}\,\frac{p^l-\xi k^l}{(\v p-\xi\v k)^2}\,e^{i\v k\cdot\v X+i\v p\cdot\v Y}=-\frac{X^l+Y^l}{\v X^2-\xi\v Y^2}\,\frac{X^m+\xi Y^m}{(\v X+\xi\v Y)^2}\,,\label{eq:FT_i5}\\
&\int\frac{d^2k}{2\pi}\frac{d^2p}{2\pi}\,\frac{1}{\xi(1-\xi)\v k^2+\v p^2}\,e^{i\v k\cdot\v X+i\v p\cdot\v Z}=\frac{1}{\v X^2+\xi(1-\xi)\v Z^2}\,,\label{eq:FT_i7}\\
&\int\frac{d^2k}{2\pi}\,\frac{k^i}{\v k^2}\,\ln\frac{\v k^2}{\mu^2}\,e^{i\v k\cdot\v X}=\frac{iX^i}{\v X^2}\Big[-2\gamma_E-\ln\frac{\v X^2\mu^2}{4}\Big]\,.\label{eq:FT_log}
\end{align}
The integrals~\eqref{eq:FT_i3}--\eqref{eq:FT_i5}, whose denominators vanish inside the integration region, are understood as principal values. Eq.~\eqref{eq:FT_K} gives the Weizs\"acker--Williams field: in the conventions used here $\int_Xe^{i\v k\cdot\v X}\,X^i/\v X^2=2\pi i\,k^i/\v k^2$. Eq.~\eqref{eq:FT_log} follows from the angular integral $\int_0^{2\pi}d\theta\cos\theta\,e^{ikX\cos\theta}=2\pi iJ_1(kX)$ and $\int_0^\infty dt\,J_1(t)\ln t=\ln2-\gamma_E$; it relates the $\MS$ logarithm in momentum space to the combination
\begin{align}\label{eq:c0}
2\gamma_E+\ln\frac{\v X^2\mu^2_{\MS}}{4}=\ln\frac{\v X^2\mu^2_{\MS}}{c_0^2}\,,\qquad c_0=2e^{-\gamma_E}\,,
\end{align}
that appears in coordinate space.

% ======================================================================
\section{Prescriptions and distributions}\label{app:distributions}
% ======================================================================

We collect here the conventions used for the singular longitudinal and transverse integrals.

\paragraph{Principal values.} Energy denominators that vanish inside the region of integration, such as $p^+\v k^2-k^+\v p^2$ in momentum space or $p^+(\v x-\v u)^2-k^+(\v x-\v z)^2$ in coordinate space, arise from transitions between states of equal free energy and are understood as Cauchy principal values. Correspondingly, logarithms that arise from integrals over such denominators are logarithms of absolute values, $\log|\cdot|$. This is the standing-wave dressing of the final state; the out-state dressing differs by the finite term $\Delta N$ of Eq.~\eqref{eq:DeltaN} (Section~\ref{sec:fin_pv}).

\paragraph{Dimensional regularisation.} Transverse loop integrals are computed in $d=2-\epsilon$ dimensions, $\int\frac{d^2p}{(2\pi)^2}\to\mu^{\epsilon}\int\frac{d^dp}{(2\pi)^d}$, in the $\MS$ scheme with scale $\mu_{\MS}$. Scaleless integrals vanish, $\int d^dp/\v p^2=0$. Coordinate-space logarithms are written with the combination~\eqref{eq:c0}.

\paragraph{Plus prescriptions in $p^+$.} For a function $f(p^+)$ regular at $p^+=0$,
\begin{align}\label{eq:plusp}
\int_0^{P}dp^+\Big[\frac{1}{p^+}\Big]_+f(p^+)\equiv\int_0^{P}dp^+\,\frac{f(p^+)-f(0)}{p^+}\,.
\end{align}
The subtracted term, integrated from $\Lambda$, is $f(0)\log(P/\Lambda)$; when $f(0)$ is the JIMWLK (target-side) structure it is part of the JIMWLK term.

\paragraph{Plus prescriptions in the momentum fraction.} With $\xi=k^+/(p^++k^+)$ the momentum fraction of the observed gluon in its parent, and for a function $F(\xi)$,
\begin{align}\label{eq:plusxi}
\int_0^1d\xi\,\frac{\xi}{(1-\xi)_+}F(\xi)&\equiv\int_0^1d\xi\,\frac{\xi F(\xi)-F(1)}{1-\xi}\,,\notag\\
\int_0^1d\xi\,\Big[\frac{1-\xi}{\xi}\Big]_+F(\xi)&\equiv\int_0^1d\xi\,\frac{1-\xi}{\xi}\big[F(\xi)-F(0)\big]\,.
\end{align}
The first is the standard prescription at $\xi\to1$ (integrated gluon soft); the second subtracts the pole at $\xi\to0$ (observed gluon soft relative to its parent), which belongs to the projectile (BFKL) evolution. The combination of the two subtractions coincides with the ``double plus'' prescription of~\cite{Kovner:2023vsy}: with $N(\xi)=\xi^2+(1-\xi)^2+\xi^2(1-\xi)^2$, $N(0)=N(1)=1$,
\begin{align}
\int_0^1\frac{d\xi}{\xi_+(1-\xi)_+}\,N(\xi)\,F(\xi)&\equiv\int_0^1d\xi\Big[\frac{N(\xi)F(\xi)}{\xi(1-\xi)}-\frac{F(0)}{\xi}-\frac{F(1)}{1-\xi}\Big]\,,\notag\\
\frac{N(\xi)}{\xi(1-\xi)}&=\xi(1-\xi)+\frac{\xi}{1-\xi}+\frac{1-\xi}{\xi}\,,
\end{align}
and $\int_0^1\frac{d\xi}{\xi_+(1-\xi)_+}N(\xi)=-\frac{11}{6}$.

\paragraph{Integrals over $p^+$ up to $\vee\to\infty$.} Some terms have no dependence on $p^+$ other than an oscillating phase. Their integral converges for $\vee\to\infty$ in the sense of distributions,
\begin{align}\label{eq:distrib}
\int_0^\infty dp^+\,e^{-iap^+}=\pi\,\delta(a)-i\,\mathcal P\frac1a\,.
\end{align}

\paragraph{Recoil function.} The recoil of the observed gluon in the rows with $\mf B_3$ and $\mf C_1$ produces, after the longitudinal integration, the function
\begin{align}\label{eq:Erecoil}
\mathcal E(s)=\int_0^1\frac{dt}{t}\big(e^{-ist}-1\big)\,,
\end{align}
which vanishes at $s=0$ and grows logarithmically for $|s|\to\infty$, $\mathcal E(s)=-\log|s|-\gamma_E-\frac{i\pi}{2}{\rm sign}(s)+O(1/s)$.

\paragraph{Longitudinal integrals.} The following integrals, with $\delta=\Lambda/k^+\to0$, produce the constants of the one-gluon wave function at order $g^3$, Eq.~\eqref{eq:xiintegrals}:
\begin{align}
\int_\delta^{1-\delta}d\xi\Big[-2+\xi(1-\xi)+\frac{1}{\xi(1-\xi)}\Big]&=-\frac{11}{6}-2\log\delta\,,\notag\\
\int_0^1d\xi\,\big[-2+\xi(1-\xi)\big]\log\big(\xi(1-\xi)\big)&=\frac{67}{18}\,,\notag\\
\int_\delta^{1-\delta}d\xi\,\frac{\log(\xi(1-\xi))}{\xi(1-\xi)}&=-\log^2\delta-\frac{\pi^2}{3}\,,
\end{align}
up to terms that vanish for $\delta\to0$.


\bibliographystyle{JHEP}
\bibliography{refs}

\end{document}
